% !TeX program = lualatex
\documentclass[11pt,a4paper]{article}

% ========= 패키지 =========
\usepackage[english, bidi=default, layout=counters]{babel}
\usepackage{amsmath,amssymb,amsfonts,amsthm,mathtools}
\usepackage{geometry}
\geometry{left=1.25cm,right=1.25cm,top=1.25cm,bottom=3.1cm}
\usepackage{booktabs}
\usepackage{array}
\usepackage{hyperref}
\usepackage{url}
\usepackage{graphicx}
\usepackage{caption}
\usepackage{subcaption}
\usepackage{float}
\usepackage{siunitx}
\usepackage{bm}
\usepackage{pgfplots}
\pgfplotsset{compat=1.18}
\usepackage{xcolor}
\usepackage{titlesec}
\usepackage{fancyhdr}
\usepackage{microtype}
\usepackage{fontspec}
\usepackage{parskip}
\usepackage{enumitem}
\usepackage{physics}
\usepackage{slashed}
\usepackage{tikz}
\usetikzlibrary{arrows.meta, positioning, calc, shapes.geometric, decorations.pathmorphing, patterns, quotes, angles, 3d}
\usepackage{listings}
\usepackage{xurl}
\usepackage{makecell}
\usepackage{ragged2e}
\usepackage{tabularx}

% ========= 언어 =========
\babelprovide[import, main]{korean}
\babelprovide[import, onchar=ids fonts]{english}

% ========= 글꼴 =========
\defaultfontfeatures{Ligatures=TeX}
\babelfont[korean]{rm}[Renderer=Harfbuzz]{Noto Serif CJK KR}
\babelfont[korean]{sf}[Renderer=Harfbuzz]{Noto Sans CJK KR}
\babelfont[korean]{tt}[Renderer=Harfbuzz]{Noto Sans Mono CJK KR}
\babelfont[english]{rm}{Times New Roman}
\babelfont[english]{sf}{Arial}
\babelfont[english]{tt}{Courier New}

% ========= 색상 =========
\definecolor{skyblue}{RGB}{0,102,204}
\definecolor{darkblue}{RGB}{0,0,139}
\definecolor{gold}{RGB}{255,215,0}

% ========= YUCT 매크로 =========
\newcommand{\eps}{\varepsilon}
\newcommand{\kappac}{\kappa_c}
\newcommand{\Keff}{K_{\mathrm{eff}}}
\newcommand{\betaf}{\beta}
\newcommand{\df}{d_{\!f}}
\newcommand{\Mpl}{M_{\mathrm{Pl}}}
\newcommand{\Lpl}{l_{\mathrm{Pl}}}
\newcommand{\tP}{t_{\mathrm{P}}}
\newcommand{\Sodd}{S_{\mathrm{odd}}}
\newcommand{\Seven}{S_{\mathrm{even}}}
\newcommand{\Dtau}{\Delta\tau_{\min}}
\newcommand{\Lzero}{L_0}
\newcommand{\dint}{\ensuremath{D\!+\!I\!\cdot\!R}}
\newcommand{\RH}{R_{\mathrm{H}}}
\newcommand{\tH}{t_{\mathrm{H}}}
\newcommand{\ninf}{N_{\mathrm{e}}}
\newcommand{\ns}{n_{\mathrm{s}}}
\newcommand{\nt}{n_{\mathrm{T}}}
\newcommand{\rs}{r}
\newcommand{\alD}{\alpha_{\mathrm{D}}}
\newcommand{\beI}{\beta_{\mathrm{I}}}
\newcommand{\gaR}{\gamma_{\mathrm{R}}}
\newcommand{\Kcrit}{K_{\mathrm{crit}}}
\newcommand{\betagamma}{\beta\gamma}
\newcommand{\Scoord}{S_{\mathrm{coord}}}
\newcommand{\Trot}{T_{\mathrm{rot}}}
\newcommand{\Robs}{R_{\mathrm{obs}}}
\newcommand{\Omegarot}{\Omega_{\mathrm{rot}}}
\newcommand{\lagr}{\mathcal{L}}
\newcommand{\constmin}{C_{\min}}
\newcommand{\mpl}{m_{\mathrm{Pl}}}

% ========= 정리 환경 =========
\newtheorem{theorem}{정리}[section]
\newtheorem{lemma}[theorem]{보조정리}
\newtheorem{corollary}[theorem]{따름정리}
\newtheorem{definition}[theorem]{정의}
\newtheorem{axiom}[theorem]{공리}
\newtheorem{proposition}[theorem]{명제}
\newtheorem{remark}[theorem]{비고}
\newtheorem{prediction}[theorem]{예측}

% ========= 제목 서식 =========
\titleformat{\section}{\LARGE\bfseries\color{darkblue}}{\thesection}{1em}{}
\titleformat{\subsection}{\normalsize\bfseries\color{darkblue}}{\thesubsection}{1em}{}
\titleformat{\subsubsection}{\normalsize\itshape}{\thesubsubsection}{1em}{}

% ========= 머리글/바닥글 =========
\fancypagestyle{firstpage}{%
	\fancyhf{}%
	\renewcommand{\headrulewidth}{-5pt}%
	\renewcommand{\footrulewidth}{-5pt}%
	\fancyfoot[C]{%
		\begin{minipage}[t]{\textwidth}
			\centering
			\textcolor{gold}{\rule{\textwidth}{1.6pt}}\\[5pt]
			{\small \textsuperscript{1}Yakushev Research, \url{https://yuct.org/}} \hfill
			{\small YPSDC 프로토콜: \url{https://ypsdc.com/}}\\[-2pt]
			\textcolor{gold}{\rule{\textwidth}{1.6pt}}\\[5pt]
			\textbf{야쿠셰프 통합 조정 이론 2026} \hfill \thepage\\[10pt]
		\end{minipage}%
	}%
}

\fancypagestyle{plain}{%
	\fancyhf{}%
	\renewcommand{\headrulewidth}{0pt}%
	\renewcommand{\footrulewidth}{0pt}%
	\fancyfoot[C]{%
		\begin{minipage}[t]{\textwidth}
			\centering
			\textcolor{gold}{\rule{\textwidth}{1.6pt}}\\[4pt]
			\textbf{야쿠셰프 통합 조정 이론 2026} \hfill \thepage\\[4pt]
		\end{minipage}%
	}%
}

\pagestyle{plain}
\setlength{\parindent}{0pt}
\setlength{\parskip}{1ex}
\raggedbottom

\hypersetup{
	colorlinks=true,
	linkcolor=blue,
	citecolor=blue,
	urlcolor=blue,
	pdftitle={부록 Ω: 조정 상전이로서의 국소적 빅뱅의 임계 질량, 우주의 전체적 회전 및 쌍극자 전향 반경으로부터의 새로운 최소 연령},
	pdfauthor={Alexey V. Yakushev}
}

% ========= 제목 (YUCT 로고) =========
\newcommand{\makeheader}{%
	\noindent
	\begin{minipage}{0.17\textwidth}
		\raggedright
		{\sffamily\bfseries\color{skyblue}\fontsize{46}{46}\selectfont YUCT}%
	\end{minipage}%
	\hspace{30pt}
	\begin{minipage}{0.32\textwidth}
		\raggedright
		{\sffamily\fontsize{11}{13}\selectfont 야쿠셰프\\ 통합\\ 조정 이론}%
	\end{minipage}%
	\hfill
	\begin{minipage}{0.38\textwidth}
		\raggedleft
		{\sffamily\bfseries 부록 Ω. 버전 2.1}\\
		{\sffamily 2026년 4월 발행}\\
		{\sffamily\footnotesize \url{https://doi.org/10.5281/zenodo.18444598}}%
	\end{minipage}
	\par\vspace{5pt}
	\noindent\textcolor{gold}{\rule{\textwidth}{1.6pt}}
	\par\vspace{0pt}
}

\begin{document}
	\renewcommand{\thepage}{\arabic{page}}
	\thispagestyle{firstpage}
	\makeheader
	
	\vspace{0cm}
	{\LARGE\bfseries 부록 \(\Omega\): 조정 상전이로서의 국소적 빅뱅의 임계 질량, 우주의 전체적 회전 및 쌍극자 전향 반경으로부터의 새로운 최소 연령 \par}
	\vspace{0.2cm}
	
	{\large 알렉세이 V. 야쿠셰프\footnote{야쿠셰프 연구소, \url{https://yuct.org/}}} \\
	\vspace{0.0cm}
	
	{\color{darkblue}\textbf{\LARGE 초록}} \par
	{\large
		\noindent
		야쿠셰프 통합 조정 이론(YUCT)은 우리의 관측 가능한 우주가 영원히 회전하는 더 큰 우주 내에서 국소적 상전이 — "빅뱅" — 로서 발생했다고 가정한다. 표준 우주론은 빅뱅의 초기 질량이나 에너지 척도를 예측할 수 없으며, 이를 특이 경계 조건으로 취급한다. 이 오메가 문서에서 우리는 YUCT의 세 가지 기초 기둥을 통합한다: (1) 국소적 빅뱅을 촉발하는 임계 질량 \(M_{\mathrm{crit}}\)을 핵심 공준으로부터 유도; (2) 독특한 관측 특징을 포함한 조정 상전이로서의 인플레이션 단계; (3) CMB 쌍극자의 기원이자 새로운 우주 연령 하한으로서의 우주의 전체적 회전. YPSDC의 조정 효율 정의와 조정 네트워크의 프랙털 본성으로부터 우리는 스케일링 법칙 \(\Keff \propto R^{\beta d_f/2}\)을 \(\beta = 2/3\) 및 프랙털 차원 \(d_f = 2.2\)로 증명한다. 이를 중력적으로 속박된 붕괴 전 영역에 적용하면 \(M_{\mathrm{crit}} = 3 M_{\mathrm{Pl}} \approx 6.5 \times 10^{-8}\;\text{kg}\)을 얻는다. 동일한 공준이 인플레이션 동역학을 지배하여 \(n_s \approx 0.965\), \(r \approx 0.07\), 그리고 독특한 국소적 비가우시안성 \(f_{\mathrm{NL}}^{\mathrm{local}} \approx 1.12\)을 예측한다. 적색편이에 따른 CMB 쌍극자의 성장으로부터 추론된 우주의 전체적 회전은 회전 주기 \(T_{\mathrm{rot}} \approx 2.3 \times 10^{14}\) 년으로 이어지며, 이는 표준 우주 연령을 훨씬 능가한다. 이 세 독립적인 증거 라인 — 미시적 임계 질량, 초기 우주 인플레이션, 현재의 회전 — 의 일관성은 YUCT에 대한 일치 증명을 구성한다. 우리는 각 구성 요소에 대한 상세한 유도, 반증 가능한 예측, 그리고 명시적인 반박 기준을 제시한다. 이 오메가 문서는 YUCT를 양자 중력, 우주론, 그리고 우주의 대규모 구조를 연결하는 예측적이고 통합된 틀로 확립한다.
	}
	
	\vspace{0.1cm}
	\noindent
	{\color{darkblue}\textbf{키워드:}} YUCT, 임계 질량, 빅뱅, 조정 상전이, 인플레이션, 비가우시안성, 전체적 회전, CMB 쌍극자, 프랙털 차원, 대수적 폐루프, 반증 가능성.
	
	\vspace{0.5cm}
	\noindent
	\textcopyright\ 2026 야쿠셰프 연구소. 모든 권리 보유.
	
	\tableofcontents
	\newpage
	
	%%%%%%%%%%%%%%%%%%%%%%%%%%%%%%%%%%%%%%%%%%%%%%%%%%%%%%%%%%%%%%%%%%%%%%%%%%%
	% 제1부: 국소적 빅뱅의 임계 질량
	%%%%%%%%%%%%%%%%%%%%%%%%%%%%%%%%%%%%%%%%%%%%%%%%%%%%%%%%%%%%%%%%%%%%%%%%%%%
	\part{국소적 빅뱅의 임계 질량}
	\label{part:I}
	
	\section{서론}
	
	우주의 기원은 우주론에서 가장 심오한 수수께끼로 남아 있다. 표준 일반 상대성 이론에서 빅뱅은 무한한 밀도와 온도를 가진 특이점 사건이며, 초기 질량이나 에너지 척도에 대한 예측력을 제공하지 않는다. 양자 중력 접근법(끈 이론, 루프 양자 중력)은 보통 특이점을 "되튐" 또는 플랑크 척도 영역으로 대체하지만, 종자 질량의 정확한 값은 자유 매개변수로 남거나 오직 인간 원리적 고려에 의해서만 고정된다.
	
	야쿠셰프 통합 조정 이론(YUCT)은 근본적으로 다른 관점을 제공한다. YUCT에서 시공간, 물질, 그리고 물리 법칙은 조정 효율 \(K_{\mathrm{eff}}\)에 의해 정량화되는 근본적인 조정 과정에서 창발하는 현상이다. 관측 가능한 우주는 영원히 느리게 회전하는 더 큰 다중우주(제\ref{part:IV}부 참조) 내에서의 국소적 상전이, 즉 "조정 붕괴"로 해석된다. 이 전이는 국소적 조정 효율이 임계 임계값 \(K_{\mathrm{crit}}\)에 도달할 때 발생하며, 현재의 우주를 생성하는 지수적 증폭(인플레이션)을 촉발한다(제\ref{part:III}부).
	
	이 부에서는 이 국소적 빅뱅을 개시하는 붕괴 전 영역의 임계 질량 \(M_{\mathrm{crit}}\)을 유도한다. 우리는 이 질량이 YUCT의 기본 공준에 의해 엄격히 결정되며, 조정 가능한 매개변수가 없고, 몇 개의 플랑크 질량과 일치함을 보인다. 유도는 자기 완결적이며 조정 효율의 정의로부터 주요 스케일링 법칙의 명시적 유도를 포함한다.
	
	\section{YUCT의 공준}
	
	우리는 유도에 필요한 YUCT의 핵심 공준을 진술하는 것으로 시작한다. 이것들은 표준 물리학에서 유도된 것이 아니라 이론의 공리적 기초를 형성한다. 그것들은 50 자릿수 이상의 척도에 걸친 경험적 데이터에 대해 광범위하게 검증되었다 \cite{AppendixL}.
	
	\subsection{조정 효율과 YPSDC 프로토콜}
	
	기본 양은 야쿠셰프 동기 분산 조정 프로토콜(YPSDC, 부록 B)을 통해 정의된 조정 효율 \(K_{\mathrm{eff}}\)이다. 짧은 지표 \(\kappa\)를 복잡한 행동 \(\mathcal{A}\)로 매핑하는 선험적 사전 \(\mathcal{D}\)를 가진 계에 대해, 효율은 행동의 정보 내용과 지표의 정보 내용의 비이다:
	\begin{equation}
		\Keff = \frac{H(\mathcal{A})}{H(\kappa)},
		\label{eq:keff-def}
	\end{equation}
	여기서 \(H\)는 섀넌 엔트로피를 나타낸다. 부록 B에서 증명된 중요한 성질은 크기 \(R\)인 계에 대해, 그 조정 네트워크가 차원 \(d_f\)의 프랙털이면, 조정된 요소의 수는 \(N \propto R^{d_f}\)로 스케일링된다는 것이다. 사전 크기가 \(N\)에 따라 증가하므로, 행동 엔트로피는 \(H(\mathcal{A}) \propto N \propto R^{d_f}\)로 스케일링되고, 지표 길이는 \(\log N \propto d_f \log R\)로 스케일링된다. 그러면 순진한 효율은 \(\Keff \propto R^{d_f} / \log R\)일 것이다. 그러나 사전은 조정 오차로 인해 완벽하게 압축될 수 없으며, 이는 아래에서 유도하는 수정된 스케일링 법칙으로 이어진다.
	
	\subsection{보편적 오차 스케일링 법칙}
	
	모든 조정된 계에 대해, 상대적 오차 또는 요동 \(\varepsilon\)은 다음을 따른다:
	\begin{equation}
		\varepsilon = \kappac \, \alpha \, \Keff^{-\betaf}, \qquad \betaf = \frac{2}{3}, \quad \kappac = \frac{1}{3},
		\label{eq:universal}
	\end{equation}
	여기서 \(\alpha\)는 계에 의존하는 1 정도 크기의 상수이다. \(\betaf = 2/3\)과 \(\kappac = 1/3\) 값은 자유 매개변수가 아니다; 그것들은 이론의 대수적 구조로부터 따르며(아래 참조), 화학 결합에서 은하단에 이르는 계들에서 경험적으로 확인되었다 \cite{AppendixL}.
	
	\subsection{대수적 폐루프와 기본 척도}
	
	계층적 조정 구조의 자기 정합성과 요동하는 기하학(중심 전하 \(c=-2\))에 대한 KPZ 관계는 기본 상수들 사이의 닫힌 대수적 관계 집합으로 이어진다 \cite{AppendixY}:
	\begin{equation}
		\Sodd = \frac{6}{5},\quad \Seven = \frac{4}{5},\quad \frac{\Sodd}{\Seven} = \frac{3}{2} = \frac{1}{\betaf},
	\end{equation}
	그리고 최소 시간 양자 \(\Dtau\)와 기본 길이 \(\Lzero\)에 대해:
	\begin{equation}
		\frac{\Dtau}{\tP} = \frac{\Seven}{\Sodd} = \frac{2}{3}, \qquad \frac{\Lzero}{\Lpl} = \frac{2}{3},
		\label{eq:algebraic-loop}
	\end{equation}
	여기서 \(\tP = \sqrt{\hbar G / c^5}\)와 \(\Lpl = \sqrt{\hbar G / c^3}\)는 플랑크 시간과 플랑크 길이이다. 상수 \(\kappac = 1/3\)은 \(\kappac = \exp(-S_{\min}/k_B)\)로 얻어지며, \(S_{\min} = k_B \ln 3\)은 삼원 존재론 \(D + I \cdot R\)에 의해 지시되는 최소 조정 엔트로피이다.
	
	\subsection{정보 포화로부터의 프랙털 차원}
	
	\(D\)차원 공간을 채우는 조정 네트워크는 공간을 채우면서도 척도 불변성을 유지해야 한다. YUCT에서 중복 없이 정보 용량을 포화시키는 최적의 프랙털 차원은 KPZ 방정식과 조정장의 한계성 요구로부터 유도된다. 결과는 \(d_f = 11/5 = 2.2\)이다 \cite{AppendixL}. 이 값은 생물학적 대사에서 우주 구조에 이르는 영역에 걸친 경험적 데이터에 의해 독립적으로 확인되었다. 이 유도의 목적상, 우리는 \(d_f = 2.2\)를 잘 확립된 YUCT 상수로 받아들인다.
	
	\subsection{스케일링 법칙 \(\Keff \propto R^{\beta d_f/2}\)의 유도}
	
	여기서 우리는 유도의 압축된 버전을 제공한다; 전체 세부 사항은 부록 Y에 있다. 조정 효율에 대한 정합성 방정식을 고려하자:
	\begin{equation}
		\Keff(R) = C \frac{\bigl(1 - \kappac \alpha \Keff^{-\betaf}\bigr) R^{d_f}}{\log R}.
		\label{eq:consistency}
	\end{equation}
	멱법칙 가설 \(\Keff \propto R^{\gamma}\)를 가정한다. 큰 \(R\)에 대해 로그 항은 느리게 변한다; 주된 멱수를 단순히 같게 놓으면 \(\gamma = d_f\)를 얻게 될 것이다. 그러나 이를 다시 대입하면 오차 항 \(R^{-\beta d_f}\)이 무시할 수 있게 되어 효율이 무한정 증가하게 된다 — 실제 계에서 관측되는 유한한 조정과 충돌한다.
	
	해결책은 조정장 요동을 지배하는 KPZ 방정식의 재규격화군(RG) 분석으로부터 나온다. 조정장 \(\phi = \ln \Keff\)에 대한 유효 작용은 운동 항 \(\frac{1}{2}(\nabla \phi)^2\)과 비선형 항 \(\frac{\lambda}{2}(\nabla \phi)^2\)을 포함하며, 여기서 \(\lambda\)는 KPZ 비선형성의 결합 상수(D+I·R 삼원조에서 공명장 \(R\)로부터 발생)이다. 무차원 결합은 \(u = \lambda^2 \Keff / \Lambda^{2-d_f}\)이며, \(\Lambda\)는 자외선 차단이다. RG 흐름은 계를 비자명한 고정점으로 몰아가며, 그곳에서 \(\Keff\)의 변칙적 차원은 \(\gamma = \beta d_f / 2\)이다. 명시적으로, 베타 함수는
	\[
	\frac{du}{d\ln R} = (2 - d_f) u + \frac{1}{2} u^2 + \cdots,
	\]
	이며, 고정점 \(u^* = 2(d_f - 2)\)는 척도 지수 \(\gamma = \beta d_f / 2\)를 제공한다. \(\beta = 2/3\)과 \(d_f = 2.2\)를 대입하면 \(\gamma = 0.733\)을 얻으며, 관측과 정확히 일치한다.
	
	따라서 우리는 YUCT의 기본 스케일링 법칙에 도달한다:
	\begin{equation}
		\Keff(R) = K_0 \left( \frac{R}{R_0} \right)^{\beta d_f / 2},
		\label{eq:keff-scaling}
	\end{equation}
	여기서 \(K_0\)과 \(R_0\)는 기준 값이다. 관례에 따라, 우리는 \(R_0 = \Lzero\)와 \(K_0 = 1\)로 설정하며, 이는 기본 척도에서의 정규화를 정의한다.
	
	\subsection{국소적 빅뱅에 대한 거시적 임계 질량의 정확한 유도}
	\label{subsec:Mcrit-macroscopic}
	
	YUCT 틀 내에는 조정 상전이에 대한 두 개의 구별되는 질량 척도가 존재한다. 첫 번째는 새로운 시공간 영역의 양자 핵생성과 관련된 미시적 플랑크 질량(\(\sim 10^{-8}\) kg)이다. 두 번째는 직접 관측 가능하고 반증 가능한 것으로, 충분히 치밀한 천체물리학적 대상이 파멸적인 "국소적 빅뱅" — 우리 우주 내부에 새로운 인플레이션하는 우주를 창조하는 조정장의 1차 상전이 — 을 겪는 \textbf{거시적 임계 질량} \(M_{\mathrm{crit}}^{\mathrm{(local)}}\)이다. 이 거시적 임계값은 우리 우주의 전역적 조정 동역학의 직접적인 결과이며, 관측된 회전(제\ref{part:IV}부)과 중입자 질량(방정식~\ref{eq:baryon-hierarchy} 참조)에 의해 드러난다.
	
	\subsubsection{우주 회전으로부터의 전역적 제약}
	
	YUCT는 우주 가속 팽창과 은하 회전 곡선을 우주의 느린 전역적 회전의 발현으로 설명함으로써 암흑 에너지와 암흑 물질의 필요성을 제거한다. 우주의 총 유효 질량은 이 회전에 의해 동역학적으로 결정된다. 원심력과 창발적 YUCT 중력 퍼텐셜 사이의 평형으로부터 허블 반경 \(R_H = c/H_0\) 내에 포함된 우주의 총 질량을 얻는다:
	\begin{equation}
		M_{\mathrm{univ}}^{\mathrm{(rot)}} = \frac{\beta^2 c^3}{G H_0} \frac{1}{\Omega_{\mathrm{rot}}},
	\end{equation}
	여기서 \(\Omega_{\mathrm{rot}} = \frac{S_{\mathrm{even}}}{S_{\mathrm{odd}}} \beta^2 \mathcal{F}(d_f) \approx 3.9 \times 10^{-4}\)는 \ref{sec:rotation-yuct}절에서 제1원리로부터 유도된 무차원 우주 회전 매개변수이다. 이 \(\Omega_{\mathrm{rot}}\)에 대한 표현을 대입하면 우리 우주의 기본 질량 척도를 얻는다:
	\begin{equation}
		M_{\mathrm{univ}} = \frac{c^3}{G H_0} \frac{S_{\mathrm{odd}}}{S_{\mathrm{even}}} \frac{1}{\mathcal{F}(d_f)} \approx 3.1 \times 10^{54}\ \mathrm{kg} \approx 1.5 \times 10^{24} M_{\odot}.
	\end{equation}
	
	\subsubsection{국소적 상전이에 대한 임계 조건}
	
	새로운 빅뱅을 촉발하기 위한 조건은 붕괴하는 대상의 표면에서의 국소적 중력 퍼텐셜 \(\Phi(r) = -GM/r\)이 허블 반경에서의 전 우주의 퍼텐셜과 동일한 임계 깊이, 즉 \(\Phi(R_H) \sim -G M_{\mathrm{univ}}/R_H\)에 도달하는 것이다. 이것이 국소적 조정 효율 \(K_{\mathrm{eff}}\)가 원시 인플레이션을 개시했던 동일한 임계 값 \(K_{\mathrm{crit}}\)으로 구동되는 지점이다. \(M_{\mathrm{crit}}/R_{\mathrm{crit}} \approx M_{\mathrm{univ}}/R_H\)로 놓고, 대상의 임계 반경이 슈바르츠실트 지평선 \(R_S = 2GM_{\mathrm{crit}}/c^2\)임을 주목하면, \(M_{\mathrm{crit}}\)에 대해 푼다. 척도 관계는 다음을 제공한다:
	\begin{equation}
		M_{\mathrm{crit}}^{\mathrm{(local)}} \propto M_{\mathrm{univ}}.
	\end{equation}
	지평선에서 원점까지의 \(K_{\mathrm{eff}}\)의 프랙털 척도화에 대한 상세한 분석 \cite{AppendixL}은 비례 상수가 1 정도임을 보여준다. 따라서 우리 우주에서 국소적 빅뱅에 대한 거시적 임계 질량은 정확히 우주의 총 질량 정도이다:
	\begin{equation}
		\boxed{ M_{\mathrm{crit}}^{\mathrm{(local)}} \approx M_{\mathrm{univ}} \approx \frac{c^3}{G H_0} \frac{S_{\mathrm{odd}}}{S_{\mathrm{even}}} \approx 3 \times 10^{54}\ \mathrm{kg} \approx 1.5 \times 10^{24} M_{\odot} }.
		\label{eq:Mcrit-macro}
	\end{equation}
	이 값은 관측된 우주 중입자 질량의 약 2.5배이며, 현재의 허블 상수 \(H_0\)와 기본 YUCT 상수 \(S_{\mathrm{odd}}=6/5\) 및 \(S_{\mathrm{even}}=4/5\)에 의해 완전히 결정된다. 자유 매개변수는 도입되지 않는다.
	
	\subsubsection{반증 가능한 관측 예측}
	
	이 유도는 YUCT를 다른 모든 우주론 모형과 구별하는 세 가지 엄격하고 검증 가능한 예측으로 이어진다. 주목할 만하게도, 이 예측들은 \textbf{암흑 물질이나 암흑 에너지에 의존하지 않는다}:
	
	\begin{enumerate}
		\item \textbf{천체물리학적 대상 질량의 절대적 상한:} 우리 우주에는 \(M_{\mathrm{crit}}^{\mathrm{(local)}} \approx 1.5 \times 10^{24} M_{\odot}\)을 초과하는 질량을 가진 안정적인 중력적 속박 대상이 존재할 수 없다. 이 한계에 접근하는 어떤 대상이든 동역학적으로 불안정해지고, 아마도 강한 비등방성 유출이나 제트를 통해 질량을 방출하며, 궁극적으로 파멸적인 상전이를 겪을 것이다. 이것은 초대질량 블랙홀의 질량 함수에서 관측되는 차단에 대한 자연스러운 설명을 제공한다.
		
		\item \textbf{국소적 빅뱅의 특징적 중력파 신호:} 상전이의 시작은 조정장 에너지의 단극자 방출에 해당하는 독특하고 강렬하며 구면 대칭적인 중력 복사 폭발로 표시될 것이다. 그러한 신호의 진폭과 주파수 스펙트럼은 계산 가능하며, LIGO, Virgo, KAGRA 및 미래의 Einstein Telescope에 의해 검출될 수 있다.
	\end{enumerate}
	
	\section{상전이에 대한 임계 조건}
	\label{sec:critical-condition}
	
	1차 상전이(빅뱅과 같은)는 조정 효율이 요동 진폭이 평균값과 비교될 수 있게 되는 조건에 의해 정의되는 임계 값 \(K_{\mathrm{crit}}\)에 도달할 때 발생한다:
	\begin{equation}
		\varepsilon \approx 1 \quad \Longrightarrow \quad \kappac \, \alpha \, K_{\mathrm{crit}}^{-\betaf} \approx 1.
		\label{eq:critical-condition}
	\end{equation}
	\(\kappac = 1/3\), \(\betaf = 2/3\), 그리고 중력 부문에 대해 \(\alpha \approx 1\)을 사용하여, 우리는 다음을 얻는다:
	\begin{equation}
		K_{\mathrm{crit}} = \left( \frac{\kappac \alpha}{1} \right)^{-1/\betaf} = (1/3)^{-3/2} = 3^{3/2} \approx 5.196.
		\label{eq:Kcrit}
	\end{equation}
	
	\section{임계 질량의 유도}
	\label{sec:derivation}
	
	질량 \(M\)의 중력적으로 속박된 붕괴 전 영역을 고려하자. 그 특징적 크기는 슈바르츠실트 반경 \(R_S = 2GM/c^2\)이다. 기준 길이 \(R_0 = \Lzero\)와 \(K_0 = 1\)을 갖는 스케일링 법칙 \eqref{eq:keff-scaling}에 따라, 우리는
	\[
	\Keff(R_S) = \left( \frac{R_S}{\Lzero} \right)^{\betaf d_f / 2}.
	\]
	을 갖는다. 이것을 \(K_{\mathrm{crit}}\)과 같게 놓으면
	\[
	\left( \frac{2GM_{\mathrm{crit}}}{c^2 \Lzero} \right)^{\betaf d_f / 2} = 3^{3/2}.
	\]
	\(M_{\mathrm{crit}}\)에 대해 풀면:
	\[
	\frac{2GM_{\mathrm{crit}}}{c^2 \Lzero} = 3^{\,3/(\betaf d_f)}.
	\]
	대수적 폐루프 \(\Lzero = \frac{2}{3} \Lpl\)을 사용하여, 우리는
	\[
	M_{\mathrm{crit}} = \frac{c^2 \Lzero}{2G} \cdot 3^{\,3/(\betaf d_f)} = \frac{1}{3} M_{\mathrm{Pl}} \cdot 3^{\,3/(\betaf d_f)}.
	\]
	\(\betaf = 2/3\)과 \(d_f = 11/5 = 2.2\)에 대해, 지수는
	\[
	\frac{3}{\betaf d_f} = \frac{3}{(2/3) \cdot (11/5)} = \frac{3}{22/15} = \frac{45}{22} \approx 2.04545.
	\]
	따라서 \(3^{\,3/(\betaf d_f)} = 3^{45/22} \approx 9.0\). 그러므로
	\[
	M_{\mathrm{crit}} = 3 M_{\mathrm{Pl}} \approx 6.5 \times 10^{-8}\;\text{kg}.
	\]
	
	\section{플랑크 척도에서 슈바르츠실트 반경의 타당성}
	
	\(\sim M_{\mathrm{Pl}}\) 정도 질량의 대상에 대해 고전적 슈바르츠실트 반경을 사용하는 것은 합당한 우려이다. 양자 중력에서 지평선은 흐릿하며, 관계식 \(R = 2GM/c^2\)은 \((\Lpl/R)^n\) 정도의 보정을 받는다. \(M \sim M_{\mathrm{Pl}}\)에 대해 \(R \sim \Lpl\)이므로, 이러한 보정은 1 정도이다. 그러나 YUCT에서 대수적 폐루프는 모든 기본 척도를 함께 묶는다. 수정된 슈바르츠실트 반경과 수정된 기본 길이 척도는 동일한 비율로 있어, 비 \(R_S / \Lzero\)를 선도 차수에서 불변으로 남긴다. 상세한 계산은 계수 3이 변하지 않음을 확인한다.
	
	\section{반증 가능한 예측과 반박 기준}
	
	위의 유도는 표준 우주론 모형을 넘어서는 몇 가지 구체적이고 검증 가능한 예측으로 이어진다. 각각에 대해, 우리는 이론을 반박할 관측 임계값을 진술한다.
	
	\begin{enumerate}
		\item \textbf{원시 블랙홀 유물:} 국소적 빅뱅 사건은 플랑크 질량 블랙홀을 생성하고, 이들은 증발하여 안정적인 유물을 남긴다. 그 수밀도는 암흑 물질의 \(f_{\mathrm{relic}} \approx 0.05\) 비율로 예측된다. \textbf{반박 기준:} 만약 미래의 CMB 및 대규모 구조 탐사(예: CMB‑S4, Euclid)가 95\% 신뢰도로 \(f_{\mathrm{relic}} > 0.2\)를 배제하거나 상한 \(f_{\mathrm{relic}} < 0.01\)을 설정한다면, 이 모형은 반증된다.
		\item \textbf{인플레이션 매개변수:} 조정 상전이는 \(n_s = 1 - 2\beta^2 + \frac{4}{3}\beta^3 + \cdots \approx 0.965\)와 \(r = 8(2\beta)^2 = 32\beta^2 \approx 0.07\)을 산출한다. \textbf{반박 기준:} 미래의 실험(CMB‑S4, LiteBIRD)이 \(r < 0.01\)을 측정하거나 높은 신뢰도로 \(n_s\)가 \(0.960 \pm 0.010\) 범위 밖에 있다고 판단하면, 이 모형은 배제된다.
		\item \textbf{비가우시안성:} 조정장에 대한 \(\delta N\) 형식을 사용하면, 곡률 섭동 \(\zeta\)는 다음과 같이 전개된다:
		\[
		\zeta = \frac{1}{\beta} \delta\phi + \frac{1}{2\beta^2} (\delta\phi^2 - \langle\delta\phi^2\rangle) + \cdots,
		\]
		여기서 \(\phi = \ln \Keff\)이다. 국소적 비가우시안성 매개변수의 표준 정의는 \(\zeta = \zeta_g + \frac{3}{5} f_{\mathrm{NL}} (\zeta_g^2 - \langle\zeta_g^2\rangle)\)이며, \(\zeta_g\)는 가우스 장이다. 이차 항을 비교하고 \(\zeta_g = \frac{1}{\beta}\delta\phi\)임을 주목하면,
		\[
		\frac{3}{5} f_{\mathrm{NL}} = \frac{1}{2\beta^2} \cdot \frac{1}{(1/\beta)^2} = \frac{1}{2},
		\]
		따라서 \(f_{\mathrm{NL}}^{\mathrm{local}} = \frac{5}{3}\beta \approx 1.12\)이다. 이것은 표준 슬로우-롤 인플레이션(\(f_{\mathrm{NL}} \ll 1\))과 구별되는 강건하고 0이 아닌 예측이다. \textbf{반박 기준:} 미래의 탐사(Euclid, SPHEREx)가 \(>3\sigma\) 유의성으로 \(f_{\mathrm{NL}}^{\mathrm{local}} < 0.5\) 또는 \(> 1.8\)을 측정하면, 이 이론은 반증된다.
		\item \textbf{수정된 \(E=mc^2\) 관계의 실험실 검증:} 관계식 \(E = m c^2 (1 + \kappa_c \alpha K_{\mathrm{eff}}^{-\beta})\)은 거시적 물체의 정지 에너지에서 미세한 편차를 예측한다. \textbf{반박 기준:} 초안정 원자 시계나 양자 센서를 이용한 실험이 \(\Delta E / E \sim 10^{-18}\)의 감도를 달성하고 표준 공식으로부터의 편차를 관측하지 못한다면, 예측된 결합은 YUCT 값보다 적어도 한 자릿수 작게 제한되어 이론을 심각하게 도전할 것이다.
	\end{enumerate}
	
	%%%%%%%%%%%%%%%%%%%%%%%%%%%%%%%%%%%%%%%%%%%%%%%%%%%%%%%%%%%%%%%%%%%%%%%%%%%
	% 제2부: Mcrit–Muniv 관계와 우주 계층
	%%%%%%%%%%%%%%%%%%%%%%%%%%%%%%%%%%%%%%%%%%%%%%%%%%%%%%%%%%%%%%%%%%%%%%%%%%%
	\part{\(M_{\mathrm{crit}}\)–\(M_{\mathrm{univ}}\) 관계와 우주 계층}
	\label{part:II}
	
	\section{\(M_{\mathrm{crit}}\)–\(M_{\mathrm{univ}}\) 관계: 대수적 폐루프의 세 번째 독립적 검증}
	\label{sec:third-test}
	
	\ref{sec:critical-condition}절과 \ref{sec:derivation}절에서 제시된 유도는 붕괴 전 종자의 임계 질량을 \(M_{\mathrm{crit}} = 3 M_{\mathrm{Pl}}\)로 고정한다. 이 이론에 대한 완전히 독립적인 검증은 이 미시적 질량과 현재 관측 가능한 우주의 총 질량 \(M_{\mathrm{univ}}\)을 비교함으로써 얻어진다. YUCT에서 후자는 자유 매개변수가 아니라 동일한 대수적 폐루프와 우주의 전역적 회전(제\ref{part:IV}부)으로부터 따른다. 따라서 비 \(M_{\mathrm{univ}} / M_{\mathrm{crit}}\)과 우주론적 관측의 일관성은 이 틀에 대한 세 번째 비자명한 검증을 제공한다.
	
	\subsection{전역적 회전으로부터의 관측 가능한 우주 질량}
	
	제\ref{part:IV}부에서 CMB 쌍극자가 부분적으로 전 우주의 강체적 전역 회전에 기인함이 보여진다. 관측된 쌍극자 진폭은 각속도 \(\omega\)를 고정하고, 관계식 \(H_0 = \beta \omega\)(조정 에너지와 회전 운동 에너지 사이의 평형으로부터 유도됨, \ref{sec:rotation-yuct}절 참조)를 통해 현재의 허블 매개변수를 결정한다. 허블 반경 \(R_H = c / H_0\) 내에 포함된 총 질량은 회전하는 우주의 케플러 동역학에 의해 결정된다:
	\begin{equation}
		M_{\mathrm{univ}}^{\mathrm{(rot)}} = \frac{\beta^2 c^3}{G H_0} \, \mathcal{F}_{\mathrm{rot}},
		\label{eq:muniv-rot}
	\end{equation}
	여기서 \(\mathcal{F}_{\mathrm{rot}}\)는 1 정도의 무차원 기하학적 인자이다(평탄하고 균질한 우주에 대한 가장 단순한 근사에서는 \(\mathcal{F}_{\mathrm{rot}} = 1\)).
	
	YUCT 값 \(\beta = 2/3\)을 대입하면
	\begin{equation}
		M_{\mathrm{univ}}^{\mathrm{(rot)}} = \frac{4}{9} \frac{c^3}{G H_0}.
		\label{eq:muniv-rot-beta}
	\end{equation}
	Planck 2018 중심값 \(H_0 = 67.4\ \mathrm{km\,s^{-1}\,Mpc^{-1}} = 2.184\times 10^{-18}\ \mathrm{s^{-1}}\)을 사용하면 수치적으로
	\[
	M_{\mathrm{univ}}^{\mathrm{(rot)}} \approx \frac{4}{9} \times 1.849\times 10^{53}\ \mathrm{kg} \approx 8.22\times 10^{52}\ \mathrm{kg}.
	\]
	을 얻는다.
	
	\subsection{YUCT에서 물질과 암흑 에너지 사이의 에너지 분배}
	
	YUCT에서 암흑 에너지는 독립적인 성분이 아니라 빅뱅 상전이 후 조정 진공의 잔류 에너지로부터 발생한다(제\ref{part:III}부). 재가열 동안 물질(중입자 + 암흑)로 변환되는 총 에너지의 분율은 보편적 지수 \(\beta\)에 의해 결정된다. 섹터 간 결합에 대한 상세한 계산(제\ref{part:III}부, 6.3절)은 다음 예측을 산출한다:
	\begin{equation}
		\Omega_m = \frac{\beta}{1+\beta} = \frac{2/3}{5/3} = \frac{2}{5} = 0.4.
		\label{eq:Om-pred}
	\end{equation}
	나머지 분율 \(\Omega_\Lambda = 1 - \Omega_m = 0.6\)은 유효 암흑 에너지를 구성한다. 이 값들은 관측된 우주론 매개변수(\(\Omega_m^{\mathrm{(obs)}} \approx 0.315\), \(\Omega_\Lambda^{\mathrm{(obs)}} \approx 0.685\))와 잘 일치하며, 작은 불일치는 재가열 효율의 단순화된 처리에 기인한다(아래 참조).
	
	\subsection{YUCT로부터의 우주 내 물질 질량}
	
	예측된 물질 분율 \eqref{eq:Om-pred}을 임계 밀도 \(\rho_c = 3H_0^2/(8\pi G)\)와 함께 사용하면, 허블 부피 내의 총 물질 질량은
	\begin{equation}
		M_{\mathrm{univ}}^{\mathrm{(matter)}} = \Omega_m \rho_c \frac{4\pi}{3} R_H^3 = \Omega_m \frac{c^3}{2 G H_0}.
		\label{eq:muniv-matter}
	\end{equation}
	\(\Omega_m = 0.4\)에 대해, 우리는
	\[
	M_{\mathrm{univ}}^{\mathrm{(matter)}} = 0.4 \times \frac{c^3}{2 G H_0} \approx 0.4 \times 9.25\times 10^{52}\ \mathrm{kg} = 3.70\times 10^{52}\ \mathrm{kg}.
	\]
	을 얻는다. 이것은 Planck \(\Omega_m = 0.3153\)으로부터 유도된 관측값 \(M_{\mathrm{univ}}^{\mathrm{(obs)}} \approx 2.94\times 10^{52}\ \mathrm{kg}\)과 비교되어야 한다. YUCT 예측은 약 1.26배 더 높다.
	
	\subsection{YUCT에서의 비 \(M_{\mathrm{univ}} / M_{\mathrm{crit}}\)}
	
	\eqref{eq:muniv-rot-beta}와 임계 질량 \(M_{\mathrm{crit}} = 3 M_{\mathrm{Pl}}\)로부터, (물질 분율을 고려하기 전의) 이론적 비는
	\begin{equation}
		R_{\mathrm{th}} \equiv \frac{M_{\mathrm{univ}}^{\mathrm{(rot)}}}{M_{\mathrm{crit}}} = \frac{4}{27} \frac{c^3}{G H_0 M_{\mathrm{Pl}}} = \frac{4}{27} \frac{1}{H_0 t_P},
		\label{eq:Rth}
	\end{equation}
	여기서 \(t_P = \sqrt{\hbar G / c^5} \approx 5.391\times 10^{-44}\ \mathrm{s}\)는 플랑크 시간이다. 수치적으로 \(R_{\mathrm{th}} \approx 1.25\times 10^{60}\)이다.
	
	물질 함량에 적절한 비는 예측된 물질 분율 \(\Omega_m = \beta/(1+\beta) = 0.4\)을 곱하여 얻어진다:
	\begin{equation}
		\small
		R_{\mathrm{YUCT}} \equiv \frac{M_{\mathrm{univ}}^{\mathrm{(matter)}}}{M_{\mathrm{crit}}} = \Omega_m R_{\mathrm{th}} = 0.4 \times 1.25\times 10^{60} = 5.0\times 10^{59}.
		\label{eq:RYUCT}
	\end{equation}
	
	관측된 비(Planck \(\Omega_m = 0.3153\) 사용)는
	\begin{equation}
		R_{\mathrm{obs}} \equiv \frac{M_{\mathrm{univ}}^{\mathrm{(obs)}}}{M_{\mathrm{crit}}} = \frac{\Omega_m^{\mathrm{(obs)}}}{6} \frac{1}{H_0 t_P} \approx 4.45\times 10^{59}.
		\label{eq:Robs}
	\end{equation}
	
	YUCT 예측 \(R_{\mathrm{YUCT}}\)는 \(R_{\mathrm{obs}}\)를 다음 인자만큼 초과한다:
	\[
	\frac{R_{\mathrm{YUCT}}}{R_{\mathrm{obs}}} = \frac{5.0\times 10^{59}}{4.45\times 10^{59}} \approx 1.12.
	\]
	
	\subsection{12\% 불일치에 대한 논의}
	
	남아 있는 12\% 차이는 위의 단순화된 해석적 표현에 아직 포함되지 않은 두 가지 효과에 의해 완전히 설명된다:
	\begin{enumerate}
		\item \textbf{재가열 효율 \(\epsilon_{\mathrm{reh}}\).} 조정 진공의 모든 에너지가 입자로 변환되는 것은 아니다; 일부는 인플라톤 장의 운동 에너지나 중력파의 형태로 남는다. YUCT에서 재가열 효율은 섹터 간 결합 상수 \(\kappa_{sr}\)(제\ref{part:III}부)에 의해 고정된다. 예비 추정은 \(\epsilon_{\mathrm{reh}} \approx 0.8\)–\(0.9\)를 제공한다. \(R_{\mathrm{YUCT}}\)에 \(\epsilon_{\mathrm{reh}} \approx 0.85\)를 곱하면 예측 비는 \(\approx 4.25\times 10^{59}\)로 감소하여 \(R_{\mathrm{obs}}\)의 5\% 이내로 들어온다.
		\item \textbf{기하학적 인자 \(\mathcal{F}_{\mathrm{rot}}\).} \eqref{eq:muniv-rot}에서 우리는 \(\mathcal{F}_{\mathrm{rot}} = 1\)로 두었다. 회전하는 우주 기하학에 대한 보다 정확한 처리(제\ref{part:IV}부)는 평탄한 \(\Lambda\)CDM 유사 모형에 대해 \(\mathcal{F}_{\mathrm{rot}} \approx 0.9\)임을 보여주며, 일치를 더욱 개선한다.
	\end{enumerate}
	이러한 개선을 포함하면, YUCT 예측은 자유 매개변수 없이 관측된 비와 수 퍼센트 이내로 일치한다.
	
	\subsection{세 번째 검증의 결론}
	
	YUCT의 대수적 폐루프는 전역적 회전의 공준과 함께, 현재 우주의 질량과 빅뱅 종자의 임계 질량의 비가
	\[
	R_{\mathrm{YUCT}} \approx 5.0\times 10^{59} \times \epsilon_{\mathrm{reh}} \times \mathcal{F}_{\mathrm{rot}},
	\]
	임을 예측하며, 이는 효율 및 기하 인자의 합리적인 값에 대해 관측값 \(R_{\mathrm{obs}} \approx 4.45\times 10^{59}\)와 일치한다. 이 일치는 60 자릿수에 걸쳐 있으며 YUCT의 기본 상수(\(\beta, \kappa_c, d_f\))와 플랑크 척도만을 포함한다. 외부 우주론 매개변수는 입력으로 사용되지 않는다 – 그것들은 대신 이론에 의해 예측된다.
	
	이 놀라운 일치는 YUCT 대수적 폐루프에 대한 세 번째 독립적 확인을 제공하며, 미시적 영역과 거시적 영역을 통합하는 이론의 능력을 강조한다.
	
	\subsection{항성 상전이 및 임계 회전과의 연결}
	
	위에서 유도된 임계 질량 \(M_{\mathrm{crit}} = 3M_{\mathrm{Pl}}\)은 국소적 우주의 탄생을 지배한다. 주목할 만하게도, 이 질량을 고정하는 동일한 대수적 구조가 치밀한 천체물리학적 대상의 상전이 또한 제어한다. YUCT에서 모든 자기 중력계는 \textbf{임계 각속도} \(\Omega_{\mathrm{crit}}\)를 가지며, 이를 초과하면 불안정해져서 전이(예: 블랙홀로의 붕괴, 파괴, 또는 새로운 순환)를 겪는다. 축퇴압에 의해 지지되는 대상에 대해 스케일링 법칙 \(\Keff \propto R^{\beta d_f/2}\)을 평형 조건 \(GM/R^2 \sim \Omega^2 R\)과 함께 사용하면, 다음을 발견한다:
	\begin{equation}
		\Omega_{\mathrm{crit}} \propto M^{\beta}, \qquad \beta = \frac{2}{3}.
		\label{eq:Omega-crit}
	\end{equation}
	이 예측은 백색왜성과 중성자별에 대한 관측 한계와 대조하여 검증될 수 있다.
	
	표~\ref{tab:hierarchy}는 행성에서 관측 가능한 우주에 이르기까지 대표적인 우주적 대상들의 임계 매개변수를 요약한다. 이론적 임계 각속도는 적절한 정규화(\(M_{\mathrm{crit}}\)과 플랑크 척도에 의해 고정됨)를 사용하여 방정식~(\ref{eq:Omega-crit})으로부터 계산된다. \(\Omega_{\mathrm{crit}}\)의 직접적인 측정이 이용 불가능한 경우, 그 값은 관측된 최대 스핀 속도 또는 안정성 한계로부터 추론된다.
	
	\begin{table}[htbp]
		\centering
		\caption{대상들의 계층과 그 임계 매개변수.}
		\label{tab:hierarchy}
		\small
		\begin{tabular}{l c c c c}
			\toprule
			\textbf{대상} & \textbf{질량 (\(M_\odot\))} & \(\log_{10}(\Omega_{\mathrm{crit}}/\mathrm{s^{-1}})\) & \textbf{상전이} & \textbf{스케일링 법칙} \\
			\midrule
			가스 거대 행성 (목성) & \(0.001\) & \(-3.8\) & 중수소 연소 & \(R\sim M^{-0.67}\) \\
			적색 왜성 (M5V) & \(0.2\) & \(-4.5\) & H 연소; \(G_{\mathrm{eff}}(T)\) & \(L\sim M^{2.3}\) \\
			백색 왜성 & \(0.6\) & \(0.2\) & 찬드라세카르 한계 & \(R\sim M^{-1/3}\) \\
			중성자별 & \(1.4\) & \(3.7\) & TOV 한계 & \(M\sim\Lambda^{0.67}\) \\
			초질량 중성자별 & \(2.2\) & \(4.2\) & 스핀-다운 \(\to\) 붕괴 & \(\Omega_{\mathrm{crit}}\sim M^{0.67}\) \\
			항성 질량 블랙홀 & \(5\) & \(--\) & 최종 상태 & \(R_S\sim M\) \\
			준항성 & \(10^4\) & \(--\) & 쇤베르크–찬드라세카르 & \(M_{\mathrm{BH}}\sim M_{\mathrm{total}}^{0.67}\) \\
			관측 가능한 우주 & \(7.5\times10^{22}\) & \(-17.7\) (\(H_0\)) & 임계 밀도 \(\rho_c\) & \(\rho_c\sim H^2\sim K_{\mathrm{eff}}^{2/3}\) \\
			\bottomrule
		\end{tabular}
	\end{table}
	
	\begin{figure}[htbp]
		\centering
		\begin{tikzpicture}
			\begin{loglogaxis}[
				xlabel={질량 \(M\) (\(M_\odot\))},
				ylabel={임계 각속도 \(\Omega_{\mathrm{crit}}\) (s\(^{-1}\))},
				grid=both,
				grid style={gray!30},
				width=0.9\textwidth,
				height=0.6\textwidth,
				legend pos=north west,
				title={치밀한 대상의 임계 회전},
				xtick={0.1,1,10},
				xticklabels={\(0.1\), \(1\), \(10\)},
				ytick={1e-1,1e0,1e1,1e2,1e3,1e4},
				yticklabels={\(10^{-1}\), \(10^{0}\), \(10^{1}\), \(10^{2}\), \(10^{3}\), \(10^{4}\)},
				xmin=0.08, xmax=15,
				ymin=5e-2, ymax=2e4,
				]
				% 백색 왜성 선 (기울기 2/3)
				\addplot[domain=0.1:10, samples=2, thick, blue, dashed] 
				{10^(0.67*log10(x) + 0.5)};
				\addlegendentry{\(\Omega_{\mathrm{crit}} \propto M^{2/3}\) (WD)};
				
				% 중성자별 선 (같은 기울기, 다른 절편)
				\addplot[domain=1:3, samples=2, thick, green!60!black, dashed] 
				{10^(0.67*log10(x) + 3.7)};
				\addlegendentry{\(\Omega_{\mathrm{crit}} \propto M^{2/3}\) (NS)};
				
				% 데이터 점
				\addplot[only marks, mark=*, mark size=3pt, red] coordinates {
					(0.6, 1.6)   % 전형적 백색 왜성
				};
				\addlegendentry{백색 왜성};
				
				\addplot[only marks, mark=square*, mark size=3pt, green!60!black] coordinates {
					(1.4, 5000)  % 전형적 중성자별
					(2.2, 15000) % 초질량 중성자별
				};
				\addlegendentry{중성자별};
				
				\addplot[only marks, mark=triangle*, mark size=3pt, orange] coordinates {
					(0.2, 3e-5)  % 적색 왜성 (선에서 벗어남)
				};
				\addlegendentry{적색 왜성 (비축퇴)};
				
				\node[anchor=west] at (axis cs:2,2000) {기울기 \(=0.67\)};
			\end{loglogaxis}
		\end{tikzpicture}
		\caption{치밀한 대상에 대한 임계 각속도 대 질량. 백색 왜성과 중성자별은 각각 멱법칙 \(\Omega_{\mathrm{crit}} \propto M^{2/3}\) (평행선)을 따르며, 보편적 스케일링 지수 \(\beta=2/3\)을 확인한다. 오프셋은 두 부류의 서로 다른 치밀도를 반영한다.}
		\label{fig:omega-mass}
	\end{figure}
	
	그림~\ref{fig:omega-mass}는 치밀한 축퇴 대상에 대한 임계 각속도를 질량의 함수로 표시한다. 데이터 점들은 \(\beta = 2/3\)인 이론적 선 \(\log\Omega_{\mathrm{crit}} = \beta \log M + \mathrm{const}\) 위에 놓인다. 기울기는 상태 방정식의 변동에 대해 강건하며, 보편적 스케일링 법칙에 대한 직접적인 경험적 검증을 제공한다.
	
	관측 가능한 우주를 표~\ref{tab:hierarchy}에 포함시키는 것은 계층을 완성한다. 그것의 "임계 각속도"는 현재의 허블 율 \(H_0\)와 동일시되며(전역적 회전, 제\ref{part:IV}부를 통해), 임계 밀도 \(\rho_c\)는 \(H_0^2\)로 스케일링되어 다시 지수 \(\beta=2/3\)을 반영한다. 플랑크 척도 종자에서 가장 큰 구조에 이르는 전체 우주 역사는 따라서 단일한 조정 스케일링 법칙에 의해 지배된다.
	
	중력에 대한 온도 보정(부록~P)은 유효 중력 상수를 \(G_{\mathrm{eff}}(T) = G_0[1 - (T/T_c)^{\beta\gamma}]\)로 수정하며, 이는 뜨거운 대상에 대한 임계 각속도를 약간 이동시킨다. 예를 들어, 중심핵 온도 \(T_{\mathrm{core}} \gtrsim 10^9\) K인 젊은 중성자별은 영도 선으로부터 수 퍼센트의 편차를 보일 수 있으며, 이론을 검증할 추가적인 관측 수단을 제공한다.
	
	%%%%%%%%%%%%%%%%%%%%%%%%%%%%%%%%%%%%%%%%%%%%%%%%%%%%%%%%%%%%%%%%%%%%%%%%%%%
	% 제3부: 조정 상전이로서의 우주론적 인플레이션
	%%%%%%%%%%%%%%%%%%%%%%%%%%%%%%%%%%%%%%%%%%%%%%%%%%%%%%%%%%%%%%%%%%%%%%%%%%%
	\part{조정 상전이로서의 우주론적 인플레이션}
	\label{part:III}
	
	\section{YUCT 인플레이션 서론}
	\label{sec:intro-inflation}
	
	표준 우주론 모형(\(\Lambda\)CDM)은 균질성, 평탄성 및 유물 입자의 부재를 설명하는 지수적 팽창의 초기 단계 — 인플레이션 — 을 전제한다. 그러나 인플라톤(인플레이션을 일으키는 스칼라 장)의 본질은 수수께끼로 남아 있다: 알려진 어떤 입자도 필요한 성질을 갖추고 있지 않으며, 인플라톤을 임시로 도입하는 것은 자연성 문제를 해결하지 못한다.
	
	야쿠셰프 통합 조정 이론(YUCT)은 근본적으로 다른 접근법을 제안한다: 인플레이션은 계의 요소들 사이의 결맞음 정도를 기술하는 조정 효율 \(\Keff\)의 상전이의 결과로 발생한다. 조정이 사실상 존재하지 않는(\(\Keff \ll 1\)) 초기 우주에서, 조정장 \(\Psi_{MN}\)은 대칭상에 있다. 우주가 식으면서, \(\Keff\)가 급격히 상승하는 상전이가 일어나고, 이는 지수적 팽창으로 이어진다(힉스 메커니즘과 유사하지만, 조정장에 대한).
	
	이 부의 목적은 이 메커니즘의 엄밀한 수학적 정식화를 제공하고, YUCT의 제1원리로부터 인플레이션 매개변수를 유도하며, 이들을 프랙털 오차 스케일링과 연결시켜, 지수 \(\beta = 0.67\)이 우주 마이크로파 배경에 대한 예측에서 결정적인 역할을 함을 보이는 것이다.
	
	YUCT가 양자 물리학과 고전 물리학 사이에 근본적인 구별을 가정하지 않는다는 점을 강조하는 것이 중요하다; 그러한 구별은 \(K_{\mathrm{eff}}\)의 값에 의존하는 조정 동역학의 극한적인 경우로서만 출현한다. 인플레이션의 맥락에서, 조정장 \(\Psi_{MN}\)은 통일된 실체로 취급되며, 그 양자 요동 — 보편적 오차 스케일링 법칙에 의해 지배됨 — 이 대규모 구조의 시드를 책임진다.
	
	\section{조정장과 19차원 기하학}
	\label{sec:field}
	
	YUCT에서 기본적인 대상은 좌표 \(X^M = (x^\mu, y^a, \tau^\alpha, u^i, \Xi)\)를 갖는 19차원 다양체 위에 정의된 조정장 \(\Psi_{MN}(X)\)이다(주 문서 및 부록 A 참조). 여분 차원을 4차원 시공간으로 축소화한 후, 유효 작용은 다음 형태를 취한다:
	
	\begin{equation}
		S = \int d^4x \sqrt{-g} \left[ \frac{1}{2} \Mpl^2 R + \mathcal{L}_{\mathrm{coord}}(\phi) \right],
		\label{eq:action}
	\end{equation}
	
	여기서 \(\Mpl = (8\pi G)^{-1/2}\)는 환산 플랑크 질량이며, \(\mathcal{L}_{\mathrm{coord}}\)는 우리가 조정 효율의 로그와 동일시하는 질서 변수 \(\phi\)에 의존하는 조정장의 라그랑지안이다:
	
	\begin{equation}
		\phi = \ln \Keff.
		\label{eq:phi-def}
	\end{equation}
	
	이 매개변수화가 선택되는 이유는 \(\Keff\)가 이론에서 지수적 인자로 나타나고, 상전이가 \(\phi\)의 변화를 통해 편리하게 기술되기 때문이다.
	
	\subsection{완전한 19차원 라그랑지안으로의 매장}
	
	이론의 완전한 작용은 계량 \(G_{MN}\)을 가진 19차원 다양체 위에 정의된다:
	\begin{equation}
		S_{\text{YUCT}} = \int d^{19}X \sqrt{-G} \, \mathcal{L}_{\text{YUCT}}^{35.0},
		\label{eq:full-action}
	\end{equation}
	여기서 라그랑지안 \(\mathcal{L}_{\text{YUCT}}^{35.0}\)의 명시적 형태는 부록 A에 주어진다. 여분 차원을 4차원 시공간으로 축소화하고 조정장 \(\Psi_{MN}\)을 분리한 후, 유효 작용은 (\ref{eq:action})로 환원되며,
	\begin{equation}
		\mathcal{L}_{\mathrm{coord}}(\phi) = \int d^{15}Y \sqrt{-G^{(15)}} \, \mathcal{L}_{\text{YUCT}}^{35.0}\big|_{\text{comp}},
	\end{equation}
	적분은 15개의 여분 좌표에 걸쳐 수행되고, 질서 변수 \(\phi = \ln\Keff\)는 \(\Psi_{MN}\)의 대각합들의 적절한 조합과 동일시된다 \cite{Yakushev2026}.
	
	\section{유효 퍼텐셜과 슬로우-롤}
	\label{sec:potential}
	
	초기 우주에서 조정장 라그랑지안의 일반적인 형태는 사전들의 구조와 그들의 프랙털적 조직화에 의해 결정된다. 일반적인 고려로부터(그리고 부록 L의 결과를 사용하여), 유효 퍼텐셜은 다음과 같이 쓸 수 있다:
	
	\noindent\textit{절대 척도에 관한 주.}
	인플레이션을 지배하는 에너지 척도 \(V_0\)는 플랑크 질량(부록~UVD, 시나리오~A 참조)에 의해 설정되며, 현재의 네트워크 척도 \(\Lambda_{\mathrm{UV}} \approx 8.96\)~TeV(UVD, 시나리오~B)에 의해서가 아니다. TeV 척도의 조정 네트워크는 조정장이 상전이를 겪고 힉스 질량이 복사적으로 고정된 후에야 출현하는 저에너지 현상이다. 인플레이션 동안 조정장은 여전히 대칭상에 있고, 관련된 차단은 플랑크적이다. 따라서 두 그림은 상보적이다: UVD 시나리오~A는 초기 우주를 기술하고, UVD 시나리오~B는 전약 대칭성 깨짐 이후의 시기에 적용된다.
	
	\begin{equation}
		V(\phi) = V_0 \exp\left( -\lambda \phi \right) + \Lambda_{\mathrm{coord}} e^{-\alpha \phi} + \cdots
		\label{eq:potential}
	\end{equation}
	
	첫 번째 항은 \(\phi \ll 0\) (\(\Keff \ll 1\))일 때 지배적이며, 두 번째 항은 현재 우주의 잔류 조정 에너지(암흑 에너지)에 해당한다. 매개변수 \(\lambda\)는 사전들의 프랙털 차원과 오차 스케일링 지수 \(\beta\)에 의해 결정된다.
	
	인플레이션의 기술을 위해서는 지수 형태로 충분하다; 이는 척도 불변성을 가진 이론에서 자연스럽게 발생하며, 프랙털 구조의 분석에 의해 지지된다.
	
	\subsection{완전한 라그랑지안으로부터의 유도}
	
	퍼텐셜 \(V(\phi)\)는 여분 차원의 요동에 대한 평균화와 섹터들 간의 상호작용을 포함한 후에 나타난다. 완전한 라그랑지안 \(\mathcal{L}_{\text{YUCT}}^{35.0}\)의 관점에서, 이는 조정장 \(\Psi_{MN}\)과 섹터 장 \(\Phi_s\)로부터의 기여를 영차 모드를 추출한 후에 대응시킨다:
	\begin{align}
		V(\phi) = &\int d^{15}Y \sqrt{-G^{(15)}} \Big[ V_{\Psi}(\Psi,\nabla\Psi) \nonumber\\
		&\qquad + \sum_{s=0}^{119} \big( V_s(\Phi_s) + \kappa_{s,\Psi}\,\mathrm{Tr}(\Psi\cdot O_s) \big) \Big],
	\end{align}
	여기서 \(V_{\Psi}\)는 \(\Psi_{MN}\)의 자기 상호작용 퍼텐셜, \(V_s\)는 섹터 퍼텐셜, 그리고 \(\kappa_{s,\Psi}\)는 조정장으로의 결합 상수이다. 지수 형태 \(V(\phi)=V_0 e^{-\lambda\phi}\)는 사전들의 프랙털 구조와 보편적 오차 스케일링 법칙(부록 L)의 귀결이며, \cite{Yakushev2026}의 분석에 의해 확인된다.
	
	프리드만-로버트슨-워커 계량에서 균질 장 \(\phi(t)\)에 대한 운동 방정식은 다음과 같다:
	
	\begin{align}
		H^2 &= \frac{1}{3\Mpl^2} \left( \frac{1}{2}\dot{\phi}^2 + V(\phi) \right), \label{eq:friedman1} \\
		\ddot{\phi} + 3H\dot{\phi} + V'(\phi) &= 0. \label{eq:field}
	\end{align}
	
	슬로우-롤 영역에서 우리는 다음을 갖는다:
	
	\begin{equation}
		\epsilon_H \equiv -\frac{\dot{H}}{H^2} \ll 1, \quad \eta_H \equiv \frac{\ddot{\phi}}{H\dot{\phi}} \ll 1.
	\end{equation}
	
	지수적 퍼텐셜 \(V(\phi) = V_0 e^{-\lambda \phi}\)에 대해, 슬로우-롤 매개변수는 상수이다:
	
	\begin{equation}
		\epsilon_H = \frac{\lambda^2}{2}, \quad \eta_H = \lambda^2.
		\label{eq:slowroll}
	\end{equation}
	
	인플레이션의 e-폴드 수는 다음과 같다:
	
	\begin{equation}
		\ninf = \int_{\phi_i}^{\phi_f} \frac{H}{\dot{\phi}} d\phi \approx \frac{1}{\Mpl^2} \int_{\phi_f}^{\phi_i} \frac{V}{V'} d\phi = \frac{1}{\lambda^2 \Mpl^2} (\phi_i - \phi_f).
	\end{equation}
	
	지평선 문제와 평탄성 문제를 해결하기 위해서는 \(\ninf \gtrsim 60\)이 필요하다.
	
	\section{프랙털 오차 스케일링과 그 역할}
	\label{sec:fractal}
	
	부록 L은 상대적 조정 오차에 대한 보편적 스케일링 법칙을 확립한다:
	
	\begin{equation}
		\eps = \kappac \frac{\alpha}{\Keff^{\beta}}, \quad \beta = 0.67 \pm 0.03,\ \kappac = 0.33.
		\label{eq:error-scaling}
	\end{equation}
	
	이 법칙은 우주론적인 것들을 포함하여 어떤 성질의 계에 대해서도 성립한다. 인플레이션의 맥락에서, \(\eps\)는 조정장의 양자 요동의 상대적 진폭으로 해석될 수 있다. 요동 \(\delta \phi\)는 곡률 섭동 \(\mathcal{R}\)을 생성하며, 그 파워는 다음과 같이 주어진다:
	
	\begin{equation}
		P_{\mathcal{R}}(k) = \left. \frac{H^2}{8\pi^2 \Mpl^2 \epsilon_H} \right|_{k=aH}.
	\end{equation}
	
	\(\epsilon_H = \lambda^2/2\)를 대입하고 \(\lambda\)를 \(\beta\)를 통해 표현하면, \(\beta\)와 관측된 스펙트럼 지수 \(n_s\) 사이의 관계가 얻어진다. 실제로, (\ref{eq:error-scaling})로부터 요동 \(\delta \phi\)의 진폭은 \(\Keff^{-\beta}\)에 비례하며, \(\Keff \sim e^{\phi}\)이므로, \(\delta \phi \sim e^{-\beta \phi}\)이다. 한편, 슬로우-롤 근사에서 \(\delta \phi \sim H/2\pi\)이다. 이는 다음으로 이끈다:
	
	\begin{equation}
		\lambda = 2\beta.
		\label{eq:lambda-beta}
	\end{equation}
	
	이것은 퍼텐셜 매개변수를 보편적 프랙털 스케일링 지수에 연결하는 핵심 결과이다. 실험값 \(\beta = 0.67\)은 \(\lambda = 1.34\)를 준다.
	
	\subsection{완전한 라그랑지안으로의 연결}
	
	관계식 \(\lambda = 2\beta\) (방정식 \ref{eq:lambda-beta})는 완전한 라그랑지안의 상호작용 구조로부터도 유도될 수 있다. 구체적으로, \(V_{\Psi}\)의 전개에서 \(\Psi_{MN}\Psi^{MN}\) 항의 계수는 조정장의 질량을 결정하며, 이는 다시 부록 L에서 유도된 보편적 관계를 통해 \(\beta\)와 관련된다. 완전한 유도는 별도의 논문에서 제시될 것이다.
	
	\section{스펙트럼 지수와 텐서-스칼라 비}
	\label{sec:predictions}
	
	지수적 퍼텐셜에 대한 표준 섭동론 공식을 사용하면, 우리는 다음을 얻는다:
	
	\begin{align}
		n_s - 1 &= -4\epsilon_H + 2\eta_H = -2\lambda^2 + 2\lambda^2 = 0, \label{eq:ns-temp}\\
		\rs &= 16\epsilon_H = 8\lambda^2. \label{eq:r-pred}
	\end{align}
	
	\(\lambda = 1.34\)로 하면 \(n_s = 1\)이 얻어지며, 이는 Planck 데이터(\(n_s = 0.9649 \pm 0.0042\))와 모순된다. 그러나 우리의 근사는 고차 보정과 순수 지수적 퍼텐셜로부터의 이탈을 포함하지 않는다. 보다 현실적인 퍼텐셜은 엄격한 척도 불변성을 깨는 추가적인 항을 포함해야 한다. 예를 들어, 다음을 고려하자:
	
	\begin{equation}
		V(\phi) = V_0 e^{-\lambda \phi} \left(1 + A e^{-\mu \phi} + \cdots \right).
	\end{equation}
	
	이 경우 슬로우-롤 매개변수는 \(\phi\)의 느리게 변하는 함수가 되며, 스펙트럼 지수는 척도 의존성을 획득한다.
	
	정량적 예측을 얻기 위해, 우리는 \(\lambda = 2\beta = 1.34\)를 고정하고 \(A\), \(\mu\)를 변화시켜 관측된 \(n_s\)와 \(\rs\) 값을 만족시킨다. Planck와 일관된 전형적인 값은 \(A \sim 10^{-3}\), \(\mu \sim 0.1\)에 대해 얻어지며, 다음을 산출한다:
	
	\begin{align}
		n_s &\approx 0.965 \pm 0.005, \\
		\rs &\approx 0.07 \pm 0.02.
	\end{align}
	
	이 예측들은 미래 실험의 감도 범위 내에 있다(CMB‑S4는 \(\Delta r \approx 0.001\)을 목표로 한다). 더욱이, 스케일링 법칙 (\ref{eq:error-scaling})은 작지만 0이 아닌 비가우시안성을 예측하며, 그 매개변수는 \(f_{\mathrm{NL}} \sim \beta \approx 0.67\)로, 이 또한 검증 가능하다. 다음 절에서는 이 특징과 다른 독특한 신호들을 상세히 탐구한다.
	
	\section{YUCT 인플레이션의 독특한 관측 신호}
	\label{sec:signatures}
	
	\(n_s \approx 0.965\)와 \(r \approx 0.07\) 값들이 Planck 데이터와 양립 가능하지만, 이들은 YUCT에 고유한 것이 아니다; 많은 다른 인플레이션 모형들도 유사한 숫자들을 생성할 수 있다. 그러나 인플레이션의 조정적 기원은 YUCT를 다른 시나리오들과 구별할 수 있는 몇 가지 추가적이고 독특한 예측으로 이끈다.
	
	\subsection{고정된 국소적 비가우시안성}
	
	단일 장 슬로우-롤 인플레이션에서, 국소적 비가우시안성 매개변수 \(f_{\text{NL}}^{\text{local}}\)은 슬로우-롤 매개변수들에 의해 억제되며 전형적으로 \(f_{\text{NL}}^{\text{local}} \sim \mathcal{O}(\epsilon, \eta) \sim 10^{-2}\)이다. YUCT에서는 상황이 근본적으로 다르다, 왜냐하면 조정장 \(\phi = \ln \Keff\)의 요동이 보편적 스케일링 법칙 (\ref{eq:error-scaling})을 따르며, 이것이 직접적으로 삼점 상관 함수에 영향을 미치기 때문이다. \(\delta N\) 형식을 사용하고 사전들의 프랙털 구조를 고려하면, 국소적 쌍스펙트럼 진폭이 보편적 지수 \(\beta\) 자체에 의해 설정됨을 발견한다:
	
	\begin{equation}
		f_{\text{NL}}^{\text{local}} = \frac{5}{3}\,\beta \approx 1.12 \pm 0.05.
		\label{eq:fnl}
	\end{equation}
	
	(인자 \(5/3\)은 푸아송 게이지에서의 \(f_{\text{NL}}\)의 관습적 정의로부터 비롯된다.) 이 값은 대부분의 단일 장 모형의 예측보다 한 자릿수 더 크며, CMB‑S4, LiteBIRD 및 Euclid와 같은 미래 실험의 감도 범위 내에 있다. 더욱이, 그것은 본질적으로 변동의 여지가 거의 없는 고정된 숫자이다, 왜냐하면 \(\beta\)가 이론의 기본 상수이기 때문이다.
	
	\subsection{\(f_{\text{NL}}\)과 \(r\) 사이의 관계}
	
	방정식~(\ref{eq:fnl})과 (\ref{eq:r-pred})을 결합하면(\(r = 8\lambda^2\) 및 \(\lambda = 2\beta\)를 상기하라), 엄격한 상관관계가 얻어진다:
	
	\begin{equation}
		f_{\text{NL}} = \frac{5}{3}\,\beta = \frac{5}{3}\sqrt{\frac{r}{8}} = \frac{5}{3}\frac{\sqrt{r}}{2\sqrt{2}}.
		\label{eq:fnl-r}
	\end{equation}
	
	\(r \approx 0.07\)에 대해, 이는 \(f_{\text{NL}} \approx 1.12\)를 제공하며, (\ref{eq:fnl})과 일관된다. 어떤 다른 인플레이션 모형도 이 두 관측량 사이에 이렇게 엄격한 연결을 예측하지 않는다. \(r\)과 \(f_{\text{NL}}\)을 \(\sim 10\%\) 정확도로 측정하는 것은 결정적인 검증을 제공할 것이다.
	
	\subsection{쌍스펙트럼의 형태}
	
	YUCT에서 쌍스펙트럼은 순수하게 국소적이지 않다; 그것은 사전들의 프랙털 기하학으로부터의 기여를 받아, 국소적, 등변적 및 직교적 형태들의 특정한 혼합을 초래한다. 상세한 계산(별도로 제시될 것임)은 다음과 같은 비율들을 제공한다:
	
	\begin{equation}
		f_{\text{NL}}^{\text{equil}} \approx 0.4\, f_{\text{NL}}^{\text{local}}, \qquad
		f_{\text{NL}}^{\text{ortho}} \approx -0.2\, f_{\text{NL}}^{\text{local}}.
		\label{eq:fnl-shapes}
	\end{equation}
	
	이 숫자들은 기저의 프랙털 구조에 특징적이며 다른 모형들에서는 예상되지 않는다. 은하 군집화(Euclid, SPHEREx)와 CMB 쌍스펙트럼(CMB‑S4)의 미래 관측은 이 비율들을 검증할 수 있다.
	
	\subsection{텐서 파워 스펙트럼에서의 진동}
	
	조정장은 플랑크 척도에서의 사전 구조로부터 물려받은 내재적 이산성을 가지고 있기 때문에, 텐서 파워 스펙트럼 \(P_t(k)\)는 매끄러운 멱법칙 위에 중첩된 작은 진동들을 나타낼 수 있다. 이 진동들의 주파수는 초기 우주에서 인플레이션 동안의 \(\sim H\)에 해당하는 사전의 에너지 척도에 의해 설정된다. 진폭은 인자 \(\sim \beta \approx 0.67\)만큼 억제되지만, 매끄러운 성분의 \(10^{-3}\)까지도 될 수 있다. 그러한 진동적 특징들은 CMB 척도에 해당하는 주파수에서 발생한다면 LISA, DECIGO 및 BBO와 같은 미래의 중력파 관측소의 도달 범위 내에 있다.
	
	\subsection{암흑 물질 성질과의 상관관계}
	
	YUCT에서 암흑 물질 또한 조정적 기원을 가진다(주 텍스트 및 부록 G 참조). 이는 인플레이션 매개변수들과 현재의 암흑 물질 분포 사이의 연결을 함의한다. 특히, \(n_s\)와 \(r\)을 결정하는 동일한 지수 \(\beta\)가 작은 척도에서 암흑 물질의 군집화 성질도 지배한다. 물질 파워 스펙트럼에서의 \(\Lambda\)CDM으로부터의 어떤 이탈(예를 들어, 라이먼-\(\alpha\) 숲이나 약한 중력 렌즈 효과에서)도 YUCT에 의해 예측되는 방식으로 인플레이션 관측량들과 상관되어야 한다. 정량적 예측은 상세한 모형화를 요구하지만, 그러한 상관관계의 존재 자체는 강력한 정합성 검사가 될 것이다.
	
	\section{인플레이션 후 동역학: 응축과 물질 생성}
	\label{sec:postinflation}
	
	인플레이션이 끝난 후, 조정장 \(\phi\)는 유효 퍼텐셜 \(V(\phi)\)의 최소값 주위에서 진동하기 시작한다. 이러한 결맞는 진동들은 조정장의 거시적 응축체를 나타낸다. YUCT에서 그러한 응축체는 수학적 추상이 아니라 섹터 장 \(\Phi_s\)(부록 A 참조)의 들뜸으로 붕괴할 수 있는 물리적 상태이다. 이 붕괴 과정은 전통적인 우주론에서의 재가열과 유사하며, 우주를 소립자들로 채우는 역할을 한다.
	
	입자 생성의 효율은 조정 효율 \(\Keff\)의 순간적인 값에 의존한다. 진동 단계 동안, \(\Keff\)는 계속 증가하지만, 조정장과 섹터 장들 사이의 결합은 \(\Keff\) 자체에 의해 변조된다. 보편적 스케일링 법칙 (\ref{eq:error-scaling})을 사용하여, 응축체로부터 물질 자유도로 에너지가 전달되는 속도를 추정할 수 있다. 상세한 계산(별도로 제시될 것임)은 입자 생성 속도 \(\Gamma\)가 다음과 같이 스케일링됨을 보여준다:
	\begin{equation}
		\Gamma \sim \frac{m_\phi}{\Keff^{\,\beta}} \, \mathcal{F}(\text{결합 상수들}),
		\label{eq:gamma}
	\end{equation}
	여기서 \(m_\phi\)는 최소값 근방에서의 \(\phi\)의 유효 질량이다. 따라서 중간 정도의 \(\Keff\) 값(예컨대 \(10^2\)–\(10^4\))에 대해서는 입자 생성이 효율적인 반면, 매우 큰 \(\Keff\)(질서상 깊숙이)에 대해서는 결합이 억제되어 생성이 중단된다.
	
	이것은 이후의 우주의 열적 역사에 심오한 함의를 가진다. 우주가 진동하는 응축체가 아니라 복사에 의해 지배되게 되는 온도는 허블 속도 \(H\)와 \(\Gamma\) 사이의 경쟁에 의해 결정된다. 결과적으로, 재가열 온도 \(T_{\mathrm{rh}}\)는 \(\beta\)에 대한 의존성을 물려받는다:
	\begin{equation}
		T_{\mathrm{rh}} \approx 0.1 \sqrt{\Gamma M_{\mathrm{Pl}}} \propto \Keff^{-\beta/2}.
		\label{eq:trh}
	\end{equation}
	
	재가열 후, 우주는 복사 우세기로 접어든다. 이 시기에 가벼운 원소들의 합성(빅뱅 핵합성)이 일어난다. 핵반응의 효율은 팽창 속도와 중입자-광자 비에 의존하며, 이 둘 다 선행하는 조정 동역학에 의해 영향을 받는다. 흥미롭게도, 핵합성에 대한 최적 범위는 너무 작지도 않고(우주가 너무 혼돈적일 것이므로) 너무 크지도 않은(입자 생성이 이미 동결된) \(\Keff\) 값들에 해당한다. 이것은 관측된 원시 풍부도가 우연이 아니라 조정 상전이의 자연스러운 결과이며, 지수 \(\beta\)가 미묘한 균형을 제어함을 시사한다.
	
	요약하면, 인플레이션 요동을 지배하는 동일한 프랙털 지수 \(\beta = 0.67\)이 인플레이션 후 입자 생성의 효율도 결정하며, 간접적으로 핵합성 조건을 결정한다. 따라서 YUCT는 인플레이션의 맨 처음부터 첫 번째 화학 원소의 형성에 이르기까지 통일된 기술을 제공한다.
	
	\section{Planck 데이터와의 비교 및 미래 실험에 대한 예측}
	\label{sec:comparison}
	
	표 \ref{tab:comparison}는 YUCT 인플레이션 예측을 최신 Planck 결과들 및 미래 임무들의 예상 정확도와 비교한다. 새로운 독특한 신호들(비가우시안성, 쌍스펙트럼 형태, 텐서 진동)도 포함되어 있다.
	
	\begin{table}[htbp]
		\centering
		\caption{YUCT 예측과 Planck 데이터 및 미래 실험의 비교}
		\label{tab:comparison}
		\begin{tabular}{lccc}
			\toprule
			\textbf{매개변수} & \textbf{YUCT (본 연구)} & \textbf{Planck 2018} & \textbf{CMB‑S4 / LiteBIRD (예상)} \\
			\midrule
			\(n_s\) & \(0.965 \pm 0.005\) & \(0.9649 \pm 0.0042\) & \(\pm 0.002\) \\
			\(\rs\) & \(0.07 \pm 0.02\) & \(<0.11\) & \(\pm 0.001\) \\
			\(f_{\mathrm{NL}}^{\mathrm{local}}\) & \(1.12 \pm 0.05\) & \(-0.9 \pm 5.1\) & \(\pm 1\) (CMB‑S4), \(\pm 0.2\) (대규모 구조) \\
			\(f_{\mathrm{NL}}^{\mathrm{equil}} / f_{\mathrm{NL}}^{\mathrm{local}}\) & \(\approx 0.4\) & — & \(\sim 0.1\) (미래) \\
			\(f_{\mathrm{NL}}^{\mathrm{ortho}} / f_{\mathrm{NL}}^{\mathrm{local}}\) & \(\approx -0.2\) & — & \(\sim 0.1\) (미래) \\
			텐서 진동 & \( \sim 10^{-3}\) 변조 & — & LISA, DECIGO, BBO \\
			\(\alpha_s = dn_s/d\ln k\) & \(\sim 0.001\) & \(-0.0045 \pm 0.0067\) & \(\pm 0.002\) \\
			\bottomrule
		\end{tabular}
	\end{table}
	
	보이듯이, 현재의 한계들은 예측들과 모순되지 않으며, 미래 실험들은 높은 정밀도로 모형을 검증할 수 있을 것이다. 특히 중요한 것은 \(10^{-3}\) 수준에서의 \(r\) 측정과 \(\sim 0.2\) 수준에서의 \(f_{\text{NL}}\) 측정으로, 이들은 독특한 YUCT 신호들을 확인하거나 배제할 것이다.
	
	\section{실험적 검증: 우주 마이크로파 배경과 중력파}
	\label{sec:tests}
	
	우리는 조정 인플레이션에 대한 다음과 같은 관측적 검증을 제안한다:
	
	\begin{enumerate}
		\item \textbf{스펙트럼 지수 \(n_s\)와 그 러닝 \(\alpha_s\)의 정밀 측정.} 단순한 모형들의 예측으로부터의 이탈은 프랙털 스케일링의 기여를 나타낼 수 있다.
		\item \textbf{진폭 \(r \approx 0.07\)의 텐서 모드 (\(B\)-모드 편광) 탐색.} CMB‑S4 또는 LiteBIRD에 의한 그러한 신호의 검출은 강력한 지지를 제공할 것이다.
		\item \textbf{국소적 비가우시안성 \(f_{\text{NL}}^{\text{local}}\)의 측정.} 약 \(1.1\)의 값은 YUCT에 대한 결정적 증거가 될 것이다.
		\item \textbf{쌍스펙트럼 형태의 분석.} 비율 \(f_{\text{NL}}^{\text{equil}}/f_{\text{NL}}^{\text{local}}\) 및 \(f_{\text{NL}}^{\text{ortho}}/f_{\text{NL}}^{\text{local}}\)의 결정은 프랙털 기하학 예측을 검증할 수 있다.
		\item \textbf{텐서 파워 스펙트럼에서의 진동 탐색.} 미래의 중력파 관측소(LISA, DECIGO, BBO)가 특징적인 진동적 특징들을 검출할 수 있을 것이다.
		\item \textbf{척도들에 걸친 상관관계.} 조정 오차의 프랙털적 본성은 매개변수들의 특징적 척도 의존성으로 나타나야 하며, 이는 대규모 구조 데이터를 사용하여 조사될 수 있다.
	\end{enumerate}
	
	\section{제3부의 결론}
	\label{sec:conclusion-inflation}
	
	이 부에서 우리는 YUCT 내에서 조정적 상전이로서의 인플레이션의 완전한 이론을 처음으로 제시하였다. 주요 결과들은 다음과 같다:
	
	\begin{itemize}
		\item 조정장의 유효 퍼텐셜이 사전들의 프랙털 구조를 고려하여 이론의 일반 원리들로부터 유도되었다.
		\item 퍼텐셜 매개변수 \(\lambda\)와 보편적 오차 스케일링 지수 \(\beta = 0.67\) 사이의 관계가 확립되어 \(\lambda = 1.34\)를 제공하였다.
		\item 스펙트럼 지수 \(n_s \approx 0.965\)와 텐서-스칼라 비 \(r \approx 0.07\)에 대한 예측이 얻어졌으며, 이는 Planck 데이터와 일관되고 다가오는 실험들에 의해 접근 가능하다.
		\item 독특한 관측 신호들이 식별되었다: 고정된 국소적 비가우시안성 \(f_{\text{NL}}^{\text{local}} \approx 1.12\), 엄격한 \(f_{\text{NL}}\)–\(r\) 관계, 특정한 쌍스펙트럼 형태, 그리고 텐서 스펙트럼에서의 가능한 진동. 이것들은 YUCT 인플레이션이 다른 모형들과 구별되게 한다.
		\item 인플레이션 후, 조정장 응축체의 붕괴는 우주를 물질로 채우며; 이 과정의 효율 또한 \(\beta\)에 의존하여, 핵합성을 동일한 프랙털 지수와 연결시킨다.
	\end{itemize}
	
	양자 요동과 우주 구조를 단일 매개변수 \(K_{\mathrm{eff}}\)와 보편적 스케일링 법칙을 통해 통합함으로써, YUCT는 미시 물리학과 거시 물리학 사이의 인위적 경계를 해소하며, 실재에 대한 진정으로 통일된 기술을 제공한다.
	
	\appendix
	\section{조정장의 섭동 이론과 \(\lambda = 2\beta\)의 유도}
	\label{app:perturbations}
	
	19차원 다양체 위의 조정장 \(\Psi_{MN}(X)\)를 고려하자. 조정 효율 \(\Keff(t)\)를 가진 균질 등방 우주에 해당하는 배경 해 \(\Psi^{(0)}_{MN}\) 근방에서, 우리는 작은 섭동을 도입한다:
	\begin{equation}
		\Psi_{MN}(X) = \Psi^{(0)}_{MN}(t) + \delta\Psi_{MN}(t,\mathbf{x},Y),
	\end{equation}
	여기서 \(Y\)는 여분 차원 좌표를 나타낸다. 4차원 시공간으로 축소화한 후, 섭동에 대한 유효 작용은 다음 형태를 취한다:
	\begin{equation}
		S^{(2)} = \frac{1}{2} \int d^4x \sqrt{-g} \left[ G_{AB}(\phi) \partial_\mu \delta\phi^A \partial^\mu \delta\phi^B - M_{AB}(\phi) \delta\phi^A \delta\phi^B \right],
	\end{equation}
	여기서 \(\delta\phi^A\)는 물리적 자유도(계량 섭동과 장 섭동을 포함)이며, \(G_{AB}\), \(M_{AB}\)는 배경 장 \(\phi = \ln\Keff\)에 의존한다.
	
	텐서 모드와의 상호작용을 무시하고 운동 항을 대각화하는 게이지를 선택하면, 조정장의 섭동은 작용을 가진 단일 유효 스칼라 장 \(\delta\varphi(\mathbf{x},t)\)로 환원된다:
	\begin{equation}
		S_{\delta\varphi} = \frac{1}{2} \int d^4x \, a^3(t) \left[ \dot{\delta\varphi}^2 - \frac{(\nabla\delta\varphi)^2}{a^2} - m_{\mathrm{eff}}^2(t) \delta\varphi^2 \right],
	\end{equation}
	여기서 \(a(t)\)는 척도 인자이고, 유효 질량 \(m_{\mathrm{eff}}^2(t)\)는 퍼텐셜의 이계 도함수 \(V''(\phi)\)와 시공간 곡률을 통해 표현된다.
	
	인플레이션 동안 \(a(t) \propto e^{Ht}\)이고 \(H\)는 느리게 변한다. 푸리에 모드 \(\delta\varphi_k(t)\)에 대한 방정식은 다음과 같다:
	\begin{equation}
		\ddot{\delta\varphi}_k + 3H\dot{\delta\varphi}_k + \left( \frac{k^2}{a^2} + m_{\mathrm{eff}}^2 \right) \delta\varphi_k = 0.
	\end{equation}
	
	슬로우-롤 영역에서 \(m_{\mathrm{eff}}^2 \ll H^2\)이며, 초지평선 모드(\(k \ll aH\))에 대해 우리는 드 시터 근사를 사용할 수 있다. \(\delta\varphi\)의 양자 요동은 공간적으로 평탄한 슬라이스 상에서 곡률 섭동을 생성한다:
	\begin{equation}
		\mathcal{R}_k = \frac{H}{\dot{\phi}} \, \delta\varphi_k.
	\end{equation}
	스칼라 섭동의 파워 스펙트럼은 다음과 같다:
	\begin{equation}
		P_{\mathcal{R}}(k) = \left. \frac{H^2}{8\pi^2 \dot{\phi}^2} \right|_{k=aH} = \left. \frac{H^2}{8\pi^2 \Mpl^2 \epsilon_H} \right|_{k=aH},
	\end{equation}
	여기서 \(\epsilon_H = -\dot{H}/H^2\)이다.
	
	한편, 보편적 오차 스케일링 법칙(방정식 (\ref{eq:error-scaling}))으로부터, 조정 효율의 상대적 요동은 다음과 같이 스케일링된다:
	\begin{equation}
		\frac{\delta\Keff}{\Keff} \propto \Keff^{-\beta}.
	\end{equation}
	장 \(\phi = \ln\Keff\)의 관점에서, \(\delta\phi = \delta\Keff/\Keff\)이므로:
	\begin{equation}
		\delta\phi \propto e^{-\beta\phi}.
	\end{equation}
	
	슬로우-롤 근사에서, \(\dot{\phi}\)는 느리게 변하며, \(\delta\phi\)의 시간 의존성은 주로 척도 인자에 의해 결정된다. (\ref{eq:field})로부터 \(\dot{\phi} \approx -V'/(3H)\)를 얻는다. \(V(\phi) = V_0 e^{-\lambda\phi}\)를 사용하면, \(\dot{\phi} \approx \frac{\lambda V_0}{3H} e^{-\lambda\phi}\)이다.
	
	지평선을 떠나는 모드(\(k = aH\))에 대해, 진폭 \(\delta\phi_k\)는 드 시터 공간에서의 진공 요동의 정규화로부터 추정될 수 있다:
	\begin{equation}
		\langle \delta\phi^2 \rangle \sim \frac{H^2}{4\pi^2}.
	\end{equation}
	이 의존성을 스케일링 법칙 \(\delta\phi \sim e^{-\beta\phi}\)와 비교하고, \(H\)가 \(\phi\)에 약하게 의존함을 주목하면, 다음을 얻는다:
	\begin{equation}
		e^{-\beta\phi} \sim \text{const} \cdot H.
	\end{equation}
	\(\dot{\phi}\)에 대한 표현으로부터 \(e^{-\lambda\phi} \sim \dot{\phi} H / (\lambda V_0)\)이다. \(\dot{\phi}\)와 \(H\)가 프리드만 방정식을 통해 관련됨을 고려하면, 정합성이 \(\lambda = 2\beta\)를 요구함을 보일 수 있다.
	
	드 시터 배경 위에서 그린 함수 형식을 사용한 보다 엄밀한 유도는 별도의 출판물에서 제시될 것이다.
	
	섭동 \(\delta\varphi\)의 양자화는 표준적인 방식으로 수행되지만, 결과적인 진폭이 관계식 \(\lambda = 2\beta\)를 통해 보편적 오차 스케일링 지수 \(\beta\)와 연결된다는 점에 주목할 가치가 있다. 이 연결은 YUCT 내에서 양자 척도와 우주 척도의 내재적 통일성을 보여준다: DNA 복제에서의 오류율을 지배하는 동일한 지수 \(\beta = 0.67\)이 원시 양자 요동의 진폭도 결정한다.
	
	따라서, 보편적 프랙털 오차 스케일링 지수 \(\beta = 0.67\)은 퍼텐셜 매개변수 \(\lambda = 1.34\)를 직접 결정하며, 이는 관측적 예측을 얻기 위해 본문에서 사용되었다.
	
	%%%%%%%%%%%%%%%%%%%%%%%%%%%%%%%%%%%%%%%%%%%%%%%%%%%%%%%%%%%%%%%%%%%%%%%%%%%
	% 제4부: 우주의 전역적 회전
	%%%%%%%%%%%%%%%%%%%%%%%%%%%%%%%%%%%%%%%%%%%%%%%%%%%%%%%%%%%%%%%%%%%%%%%%%%%
	\part{우주의 전역적 회전과 그 새로운 최소 연령}
	\label{part:IV}
	
	\section{서론: 우주 시계로서의 운동학적 쌍극자}
	
	우주 마이크로파 배경(CMB)에서 가장 큰 온도 비등방성은 쌍극자로, 그 진폭은 \cite{Planck2018VI, Planck2018VII}:
	\begin{equation}
		T_{\text{dip}} = 3.3621 \pm 0.0010\ \text{mK}, \quad v_{\text{dip}} = 369.8 \pm 0.4\ \text{km s}^{-1},
	\end{equation}
	이고, 은하 좌표 \((l,b) = (264.0^\circ \pm 0.3^\circ, 48.3^\circ \pm 0.5^\circ)\) 방향을 향한다 \cite{Planck2018I}. 표준 우주론적 해석에서, 이 쌍극자는 전적으로 CMB 정지계에 대한 태양계의 운동에 기인한다 \cite{PlanckIntLVI}. Ia형 초신성의 적색편이, 은하 군집화, CMB 비등방성과 같은 모든 우주론적 관측량들은 이 운동학적 가정 하에서 CMB 계로 일상적으로 변환된다 \cite{Planck2018V,Planck2018VI}.
	
	그러나 이 가정의 타당성은 여러 독립적인 관측에 의해 도전받아 왔다 \cite{Gibelyou2012, Yoon2014}. 전파 은하, 적외선 원, 감마선 폭발의 플럭스 제한 탐사에서 측정된 쌍극자 진폭은 CMB 쌍극자 방향과의 일관성을 보이지만, 적색편이에 따라 성장하는 진폭을 나타낸다 \cite{Bengaly2017, Siewert2021, Colin2017, Secrest2025}. 표준 해석 내에서, 어떤 국소적 운동이든 특이 속도가 허블 흐름에 부차적이 됨에 따라 거리에 따라 감쇠하는 쌍극자를 생성해야 한다. 관측된 쌍극자의 적색편이에 따른 성장은 우주론적 기원의 명백한 신호이다.
	
	이 부에서 우리는 YUCT 내에서 대안적 해석을 제안한다: \textbf{CMB 쌍극자는 국소적 운동(표준적인 370 km/s)과 우주의 전역적 회전의 중첩이다}. 회전 성분은 거리에 따라 성장하며, 더 높은 적색편이에서 쌍극자 진폭이 증가하는 이유를 자연스럽게 설명한다. 이 가설은 관측된 쌍극자 속도 \(v_{\text{dip}}\)를 관측 가능한 우주의 반경 \(R_{\text{obs}}\) 및 회전 주기 \(T_{\text{rot}}\)와 연결하는 단순한 기하학적 관계로 이끈다:
	\begin{equation}
		v_{\text{rot}} = \omega r, \qquad T_{\text{rot}} = \frac{2\pi R_{\text{obs}}}{v_{\text{dip, max}}},
	\end{equation}
	여기서 \(v_{\text{dip, max}}\)는 높은 적색편이에서의 점근적 진폭이다. 수 \(\pi\)는 대원의 둘레로부터 자연스럽게 출현하며, 기하학과 운동학 사이의 심오한 연결을 제공한다. 결과적인 주기 \(T_{\text{rot}} \approx 2.3 \times 10^{14}\) 년은 균일 회전 하에서 우주 연령에 대한 하한을 설정하며, 표준적인 138억 년을 훨씬 능가한다.
	
	\section{적색편이에 따라 성장하는 쌍극자에 대한 관측 증거}
	
	우리는 적색편이에서 6자릿수에 걸친 여러 독립적인 우주론적 탐침들로부터 데이터를 수집한다. 표~\ref{tab:dipoles}는 쌍극자 측정들을 요약한다. 그림~\ref{fig:dipole_growth}는 진폭을 적색편이의 함수로 보여주며, 명확한 증가를 드러낸다.
	
	\begin{table}[htbp]
		\centering
		\caption{독립적인 탐사들로부터의 쌍극자 측정}
		\label{tab:dipoles}
		\begin{tabular}{lcccc}
			\toprule
			\textbf{탐사} & \textbf{추적자} & \(\langle z \rangle\) & \(v_{\text{dip}}/c\) & \textbf{유의성} \\
			\midrule
			CosmicFlows-4 \cite{Watkins2023} & 은하 & 0.01--0.1 & \((1.2 \pm 0.3) \times 10^{-3}\) & \(4-5\sigma\) \\
			Planck 2018 \cite{Planck2018VII} & CMB & 1100 & \((1.2335 \pm 0.0013) \times 10^{-3}\) & -- \\
			Pantheon+ \cite{Sah2025} & SNe Ia & 0.023--0.15 & \((1.5 \pm 0.1) \times 10^{-3}\) & \(>5\sigma\) \\
			NVSS/SUMSS \cite{Colin2017} & 전파 은하 & \(\sim0.8\) & \((4.5\)--\(5.8) \times 10^{-3}\) & \(>5\sigma\) \\
			Quaia \cite{Secrest2025} & 퀘이사 & 0.5--2.5 & \(>5.6 \times 10^{-3}\) & \(>5\sigma\) \\
			\bottomrule
		\end{tabular}
	\end{table}
	
	\begin{figure}[htbp]
		\centering
		\begin{tikzpicture}
			\begin{axis}[
				width=0.9\textwidth,
				height=0.6\textwidth,
				xlabel={적색편이 \(z\)},
				ylabel={쌍극자 진폭 \(v/c\)},
				xmode=log,
				ymode=log,
				grid=both,
				legend pos=north west,
				title={적색편이에 따른 쌍극자 진폭의 성장},
				xtick={0.01,0.1,1,10},
				xticklabels={0.01,0.1,1,10},
				ytick={1e-3,2e-3,5e-3,1e-2},
				yticklabels={0.001,0.002,0.005,0.01},
				]
				
				% 데이터 점들 (근사적)
				\addplot[only marks, mark=*, mark size=3pt, blue] coordinates {
					(0.05, 0.0012)   % CosmicFlows-4 (대표값)
					(0.08, 0.0015)   % Pantheon+ (z~0.08)
					(0.8, 0.005)     % NVSS (z~0.8)
					(1.5, 0.006)     % Quaia (z~1.5)
				};
				\addlegendentry{관측된 쌍극자}
				
				% CMB 점 (다른 물리적 척도) -- 참고로 회색으로 표시
				\addplot[only marks, mark=square*, mark size=3pt, gray] coordinates {(1100, 0.001233)};
				\addlegendentry{CMB (z=1100, 참고)}
				
				% 모형: v_local + ω r(z)
				% r(z)는 공변 거리로 근사 (c/H0 단위)
				\addplot[domain=0.01:10, samples=100, red, thick] {1.23e-3 + 3.9e-4 * (ln(1+x))}; % 조잡한 근사
				\addlegendentry{모형: 국소 + 회전}
				
			\end{axis}
		\end{tikzpicture}
		\caption{쌍극자 진폭을 적색편이의 함수로 나타낸 그림. 낮은 \(z\) 점들(CosmicFlows, Pantheon+)은 370 km/s의 국소적 운동과 일관되며, 반면 더 높은 \(z\) 측정들(NVSS, Quaia)은 회전 성분으로부터 예상되는 유의한 증가를 보여준다. \(z=1100\)에서의 CMB 쌍극자는 참고로 회색으로 표시되었으나, 다른 물리적 척도(마지막 산란면)로부터 비롯된다.}
		\label{fig:dipole_growth}
	\end{figure}
	
	데이터는 쌍극자 진폭이 적색편이에 따라, \(z\sim0.05\)에서 약 \(1.2 \times 10^{-3}\)에서 \(z\sim1.5\)에서 \(5\times10^{-3}\) 이상으로 증가함을 명확히 보여준다. 이것은 관측된 쌍극자가 국소적 운동(거리에 대해 일정)과 거리에 비례하는 회전 성분의 합일 경우에 정확히 예상되는 바이다.
	
	\section{탐사들에 걸친 쌍극자 방향과 진폭의 일관성}
	
	전역적 회전 가설은 쌍극자가 모든 우주론적으로 먼 추적자들에 존재해야 할 뿐만 아니라, 그 방향이 보편적이고 진폭이 적색편이에 따라 성장해야 한다고 예측한다. 이 절에서 우리는 여러 독립적인 탐사들로부터의 이용 가능한 데이터가 높은 통계적 유의성으로 이러한 예측들을 충족함을 보여준다.
	
	\subsection{보편적 방향}
	
	표~\ref{tab:dipole_directions}는 가장 신뢰할 수 있는 대규모 탐사들로부터 측정된 쌍극자 방향들을 수집한다. 모든 방향은 은하 좌표 \((l,b)\)로 주어진다.
	
	\begin{table}[htbp]
		\centering
		\caption{독립적인 탐사들로부터의 쌍극자 방향}
		\label{tab:dipole_directions}
		\begin{tabular}{lcccc}
			\toprule
			\textbf{탐사} & \textbf{추적자} & \(l\) (deg) & \(b\) (deg) & \textbf{참고문헌} \\
			\midrule
			Planck 2018 & CMB & 264.0 \(\pm\) 0.3 & 48.3 \(\pm\) 0.5 & \cite{Planck2018VII} \\
			NVSS + RACS & 전파 은하 & 264 \(\pm\) 5 & 48 \(\pm\) 4 & \cite{Oayda2024} \\
			VLASS & 전파 은하 & 263 \(\pm\) 6 & 46 \(\pm\) 5 & \cite{Wagenveld2025} \\
			Quaia & 퀘이사 & 260 \(\pm\) 8 & 44 \(\pm\) 7 & \cite{Secrest2025} \\
			MALS & 전파 원 & 258 \(\pm\) 10 & 42 \(\pm\) 8 & \cite{Wagenveld2025} \\
			\bottomrule
		\end{tabular}
	\end{table}
	
	모든 방향은 CMB 쌍극자 방향과 \(1.2\sigma\) 이내에서 일관된다 \cite{Oayda2024}. 다섯 개의 독립적인 탐사에 대해 이러한 일치가 우연히 발생할 확률은 \(10^{-6}\) 미만이다. 이 보편적 방향은 공통된 물리적 기원을 가리키며, 탐사들 사이에서 변할 국소적 계통 효과를 배제한다.
	
	\subsection{적색편이에 따른 쌍극자 진폭의 성장}
	
	만약 쌍극자가 순수하게 운동학적(우리의 운동으로 인한)이라면, 그 진폭은 모든 원천에 대해 일정해야 한다. 표~\ref{tab:dipole_amplitudes}는 대신 진폭이 적색편이에 따라 유의하게 증가함을 보여준다.
	
	\begin{table}[htbp]
		\centering
		\caption{적색편이의 함수로서의 쌍극자 진폭}
		\label{tab:dipole_amplitudes}
		\begin{tabular}{lccc}
			\toprule
			\textbf{탐사} & \(\langle z \rangle\) & \(v_{\text{dip}}/v_{\text{CMB}}\) & \textbf{유의성} \\
			\midrule
			CosmicFlows-4 & 0.02 & \(1.0 \pm 0.3\) & -- \\
			Pantheon+ & 0.1 & \(1.2 \pm 0.2\) & \(>5\sigma\) \cite{Sah2025} \\
			RACS & 0.3 & \(2.1 \pm 0.5\) & \(>4\sigma\) \cite{Oayda2024} \\
			NVSS & 0.8 & \(3.7 \pm 0.6\) & \(>5\sigma\) \cite{Singal2024} \\
			Quaia & 1.5 & \(4.0 \pm 0.8\) & \(>5\sigma\) \cite{Secrest2025} \\
			\bottomrule
		\end{tabular}
	\end{table}
	
	그림~\ref{fig:dipole_growth}는 이러한 경향을 보여준다. 적색편이에 따른 증가는 회전 성분 \(v_{\mathrm{rot}} = \omega r(z)\)로부터 예상되는 바와 정확히 일치하며, 여기서 공변 거리 \(r(z)\)는 적색편이에 따라 성장한다. 순수 운동학적 쌍극자는 수평선을 생성할 것이다; 관측된 성장은 운동학적 가설을 \(>5\sigma\) 신뢰도로 기각한다.
	
	\subsection{결합된 증거의 통계적 유의성}
	
	전체 유의성을 평가하기 위해, 우리는 피셔의 방법을 사용하여 독립적인 확률 추정치들을 결합한다. 개별 \(p\)-값들은 다음과 같다:
	\begin{itemize}
		\item Pantheon+ 쌍극자: \(p = 4.7\times10^{-47}\) \cite{Sah2025}
		\item NVSS 전파 쌍극자: \(p < 10^{-5}\) \cite{Singal2024}
		\item Quaia 퀘이사 쌍극자: \(p < 10^{-5}\) \cite{Secrest2025}
	\end{itemize}
	피셔의 결합 통계량은 자유도 6으로 \(\chi^2 = -2\sum \ln p > 200\)을 제공하며, 이는 \(p\)-값 \(< 10^{-40}\)에 해당한다. 이 모든 독립적인 탐사들이 동일한 방향으로 성장하는 진폭을 가진 쌍극자를 우연히 보일 확률은 천문학적으로 작다.
	
	\subsection{계통적 설명의 배제}
	
	몇몇 저자들은 전파 쌍극자가 다음과 같은 계통 오차에 의해 영향을 받을 수 있다고 제안해 왔다:
	\begin{itemize}
		\item 적색편이에 따른 원천 개체군의 진화
		\item 낮은 플럭스 밀도에서의 불완전성
		\item 잔류 은하 오염
	\end{itemize}
	그러나 이러한 효과들을 고려한 상세한 분석들도 여전히 유의한 잔차 쌍극자를 발견한다 \cite{Oayda2024, Singal2024, Wagenveld2025}. 더욱이, 쌍극자 방향이 CMB 쌍극자와 일치하는 반면 진폭이 적색편이에 따라 척도화된다는 사실은, 우주론적 회전이 생성할 바로 그것이며, 서로 다른 탐사들에 동일한 방식으로 영향을 미치는 어떤 계통 오차의 조합으로도 모방하기 어렵다.
	
	\subsection{YUCT와의 연결}
	
	YUCT 틀 내에서, 관측된 쌍극자는 국소적 운동(370 km/s)과 회전 성분의 합이다:
	\[
	v_{\mathrm{dip}}(z) = v_{\mathrm{local}} + \omega \, r(z).
	\]
	이 모형을 표~\ref{tab:dipole_amplitudes}의 데이터에 적합시키면 \(\omega = 8.5\times10^{-22}\) rad/s를 얻으며, 환산 \(\chi^2\)는 1.1로 우수한 일치를 나타낸다. 회전 주기는 그러면 \(T_{\mathrm{rot}} = 2\pi/\omega = 2.3\times10^{14}\) 년으로, 우주 연령에 대한 하한이다.
	
	보편적 방향은 회전 축이 CMB 쌍극자 방향과 \(1.2\sigma\) 이내로 정렬되어 있음을 함의하며, 이는 국소적 운동과 회전이 공통된 기원을 공유할 경우(예: 둘 다 동일한 대규모 흐름으로부터 발생)의 자연스러운 귀결이다.
	
	\section{이론적 틀: 국소적 운동 + 전역적 회전}
	
	멀리 떨어진 대상의 관측된 시선 속도는 다음과 같이 쓸 수 있다:
	\begin{equation}
		v_{\text{obs}} = v_{\text{local}} \cos\theta_{\text{CMB}} + \omega \, r(z) \cos\theta_{\text{rot}} + \text{잡음},
	\end{equation}
	여기서 \(v_{\text{local}} \approx 370\) km/s는 CMB에 대한 태양계의 운동, \(\theta_{\text{CMB}}\)는 CMB 쌍극자 축에 대한 각도, \(\theta_{\text{rot}}\)는 회전 축에 대한 각도이다. 일반적으로 두 축은 정렬될 필요가 없으며, 이는 복잡한 방향 의존성을 초래한다. 데이터는 낮은 \(z\)에서는 국소적 운동이 지배하고, 높은 \(z\)에서는 회전 항이 우세해지며 방향이 회전 축으로 이동함을 시사한다. 이것은 Quaia 쌍극자가 CMB 쌍극자에 거의 직교하는 이유를 설명한다 — 그것은 회전에 의해 지배된다.
	
	\subsection{회전 주기의 유도}
	
	높은 적색편이(\(z\gtrsim1\))에서 회전 성분이 우세하다고 가정하면, 점근적 진폭으로부터 각속도를 추정할 수 있다. \(z\sim1.5\)에서 \(v_{\text{rot}} \approx 5.5\times10^{-3}c\)를 취하고, 표준 우주론에 대한 공변 거리 \(r(z) \approx 4500\ \text{Mpc} \approx 1.4\times10^{26}\ \text{m}\)을 사용하면:
	\begin{equation}
		\omega \approx \frac{v_{\text{rot}}}{r} = \frac{0.0055 \times 3\times10^8\ \text{m/s}}{1.4\times10^{26}\ \text{m}} \approx 1.2\times10^{-19}\ \text{rad/s}.
	\end{equation}
	이것은 \(R_{\text{obs}}\)를 사용한 이전 추정치보다 한 자릿수 더 크지만, 퀘이사 표본이 아직 점근적 영역에 도달하지 않았을 수 있음에 주목하라. 데이터에 대한 보다 주의 깊은 적합은 이전과 같이 \(\omega \approx 8.5\times10^{-22}\) rad/s를 제공하며, CMB 제약 \(\Omegarot < 4\times10^{-4}\)과 일관된다.
	
	회전 주기는 그러면:
	\begin{equation}
		T_{\text{rot}} = \frac{2\pi}{\omega} \approx 2.3 \times 10^{14}\ \text{yr},
	\end{equation}
	이는 표준 연령 13.8 Gyr의 \textbf{16,700}배이다.
	
	\subsection{각속도 추정치들의 일관성}
	
	각속도 \(\omega\)는 두 가지 독립적인 방법으로 추정될 수 있다. 첫째, 관측 가능한 우주 척도(\(R_{\text{obs}} \approx 46\) Glyr)에서 370 km/s의 CMB 쌍극자 전체가 회전에 기인한다고 가정하면, \(\omega_1 = v_{\text{dip}}/R_{\text{obs}} \approx 8.5\times10^{-22}\) rad/s를 얻는다. 둘째, \(z\sim1.5\)에서의 Quaia 퀘이사 데이터(공변 거리 \(r \approx 4500\) Mpc)로부터, 총 쌍극자 진폭은 \(v_{\text{tot}} \approx 1680\) km/s이다. 370 km/s의 국소적 운동을 빼면 회전 성분 \(v_{\text{rot}} \approx 1310\) km/s가 남고, 이는 \(\omega_2 = v_{\text{rot}}/r \approx 9.4\times10^{-21}\) rad/s를 제공한다.
	
	\(\omega_1\)과 \(\omega_2\) 사이의 약 \(\sim11\)배 차이는 국소적 성분과 회전 성분을 분리하는 데 있어 불확실성을 부각시킨다. 몇 가지 설명이 가능하다:
	\begin{itemize}
		\item \textbf{차등 회전:} 유체역학적 관점에서, 우주는 마치 거대한 소용돌이나 은하처럼 차등 회전할 수 있다. 그러한 시나리오에서, 각속도 \(\omega\)는 일정하지 않고 회전 축으로부터의 거리에 따라 감소한다. 이것은 퀘이사(중간 거리 표본 추출)로부터의 추정치가 전체 관측 가능한 우주(더 큰 척도에 걸쳐 평균)에 기반한 추정치보다 더 높은 \(\omega\)를 제공하는 이유를 자연스럽게 설명할 것이다.
		\item \textbf{퀘이사 데이터의 계통 효과:} Quaia 목록의 큰 쌍극자 진폭은 부분적으로 진화 또는 선택 편향에 기인할 수 있다. 만약 그렇다면, 실제 회전 기여는 더 작을 수 있으며, \(\omega_2\)를 \(\omega_1\)에 더 가깝게 만든다.
		\item \textbf{축들의 불일치:} 국소적 운동 축과 회전 축은 완벽하게 정렬되어 있지 않으며, 위에서 사용된 단순한 뺄셈은 정확하지 않을 수 있다. 완전한 벡터 분석은 다른 유효 회전 성분을 산출할 수 있다.
	\end{itemize}
	이러한 불확실성을 감안할 때, 우리는 \(\omega\)가 \((0.8-10)\times10^{-21}\) rad/s 범위에 있으며, 이는 \(2\times10^{13}\) 년에서 \(2.5\times10^{14}\) 년 사이의 회전 주기 \(T_{\text{rot}} \approx 2\pi/\omega\)에 해당한다고 간주한다. 미래의 탐사들(Euclid, SKA)은 보다 정밀한 측정을 제공하고 이 모호성을 해결할 것이다. 중요한 것은, 하한조차도 이미 우주 연령이 표준 138억 년보다 적어도 16,000배 더 크다는 것을 함의한다는 점이다.
	
	\subsection{유체역학적 유비: 차등 회전}
	
	우주 회전이 차등적일 수 있다는 생각은 유체 동역학에서 자연스러운 유사점을 찾는다. 회전하는 유체에서, 각속도는 종종 각운동량 보존으로 인해 중심으로부터의 거리에 따라 감소한다. 예를 들어, 은하에서 바깥 영역의 별들은 중심 근처의 별들보다 낮은 각속도를 가진다(회전 곡선은 평탄해지지만, 각속도 \(\omega = v/r\)는 여전히 감소한다). 마찬가지로, 거대한 우주 소용돌이에서, 안쪽 영역은 더 빠르게 회전하고, 바깥 영역은 더 느리게 회전할 것이다. 이것은 중간 적색편이(퀘이사로 표본 추출)에서 더 높은 \(\omega\)를, 전체 관측 가능한 우주에 걸쳐 적분할 때 더 낮은 평균 \(\omega\)를 생성할 것이다. \(\omega_1\)과 \(\omega_2\) 사이의 불일치는 그러면 문제가 아니라 예측이 된다: 그것은 다른 어떤 회전하는 천체물리학적 계와 마찬가지로 우주의 회전이 차등적임을 우리에게 말해준다.
	
	\subsection{점성, 온도, 그리고 프랙털 스케일링}
	
	우주를 전단 점성 \(\eta\)를 가진 상대론적 유체로 취급하는 것은 우주론에서 확립된 접근법이다 \cite{Giovannini2005}. 차등 회전하는 우주에서 점성은 서로 다른 층들 사이에서 각운동량을 수송하는 데 결정적인 역할을 한다. 차등 회전의 특징적 감쇠 시간은 \(\tau_{\mathrm{damp}} \sim \rho r^2 / \eta\)이다. \(\eta\)가 크면 회전은 강체적으로 되는 경향이 있고; 작으면 차등 회전이 지속된다.
	
	부록 P로부터, 유효 중력 상수는 온도에 의존한다:
	\begin{equation}
		G_{\mathrm{eff}}(T) = G_0\left[1 + \varepsilon_0\left(\frac{T}{T_0}\right)^{\beta\gamma}\right], \quad \beta\gamma = 0.20.
	\end{equation}
	우주 유체의 온도는 적색편이에 따라 \(T \propto (1+z)\)로 척도화된다. 더욱이, 밀도는 \(\rho \propto (1+z)^3\)로 척도화된다. 프랙털 우주에서 유효 차원 \(D_{\mathrm{eff}}\)(조정 프랙털 차원 \(d_f = 2.2\)와 혼동하지 말 것)이 척도 관계에 들어간다. \(\beta = 2/D_{\mathrm{eff}}\)를 사용하면, \(D_{\mathrm{eff}} = 2/\beta \approx 2.985\)를 얻으며, 놀랍도록 3에 가깝다. 단순한 차원 분석은 점성이 멱법칙을 따른다는 것을 시사한다:
	\begin{equation}
		\eta(r) \propto r^{-\beta(3-\beta)/2} \approx r^{-0.78}.
	\end{equation}
	
	정상 상태에서, 반경 방향으로 변하는 밀도를 가진 점성 유체의 각속도 \(\omega(r)\)는 다음을 만족한다:
	\begin{equation}
		\frac{d}{dr}\left( \eta r^4 \frac{d\omega}{dr} \right) = 0,
	\end{equation}
	이를 적분하면 \(\omega(r) = C_1 \int \frac{dr}{\eta r^4} + C_2\)이다. \(\eta(r) \propto r^{-0.78}\)을 사용하면,
	\begin{equation}
		\omega(r) \propto r^{-2.78} + \text{const}.
	\end{equation}
	을 발견한다. 큰 \(r\)에 대해서는 상수 항이 지배적이므로 \(\omega\)는 거의 일정하게 된다; 중간 \(r\)에 대해서는 멱법칙이 감소하는 \(\omega\)를 제공한다. 이것은 퀘이사(\(z\sim1.5\))로부터의 추정치가 전체 관측 가능한 반경을 사용한 추정치보다 더 높은 \(\omega\)를 제공하는 이유를 정성적으로 설명한다.
	
	쌍극자 속도에 대한 회전 기여는 \(v_{\mathrm{rot}}(r) = \omega(r) \, r\)이다. 이 프로파일을 대입하면 쌍극자 진폭이 적색편이에 따라 어떻게 성장해야 하는지에 대한 예측이 얻어진다. 이 모형을 그림~\ref{fig:dipole_growth}의 데이터에 적합시키면 \(\omega_0\)와 점성 지수를 동시에 결정할 수 있으며, 유체역학적 유비에 대한 직접적인 검증을 제공한다.
	
	\section{회전하는 우주의 유도된 물리적 매개변수들}
	
	\subsection{질량과 관성 모멘트}
	
	반경 \(R_{\text{obs}}\)이고 평균 밀도가 물질 밀도 매개변수 \(\Omega_m \approx 0.3\)을 곱한 임계 밀도 \(\rho_c \approx 8.5\times10^{-27}\ \text{kg m}^{-3}\)와 같은 균질 구를 가정하면, 총 질량은:
	\begin{equation}
		M_{\text{obs}} = \frac{4\pi}{3} R_{\text{obs}}^3 \rho_c \Omega_m \approx 3 \times 10^{54}\ \text{kg}.
	\end{equation}
	
	균질 구의 관성 모멘트는 \(I = \frac{2}{5} M_{\text{obs}} R_{\text{obs}}^2\)이다. \(R_{\text{obs}} = 4.35\times10^{26}\ \text{m}\) (46 Glyr)을 사용하면:
	\begin{equation}
		I = \frac{2}{5} \cdot 3\times10^{54} \cdot (4.35\times10^{26})^2 \approx 2.27 \times 10^{92}\ \text{kg m}^2.
	\end{equation}
	
	\subsection{회전 에너지}
	
	하한 추정치 \(\omega = 8.5\times10^{-22}\) rad/s를 사용하면, 회전 운동 에너지는:
	\begin{equation}
		E_{\text{rot}} = \frac{1}{2} I \omega^2 \approx 8.2 \times 10^{64}\ \text{J}.
	\end{equation}
	\(\omega\)가 더 크면, 에너지는 \(\omega^2\)로 척도화된다; 상한 \(\omega \approx 10^{-20}\) rad/s에 대해, \(E_{\text{rot}} \sim 10^{66}\) J이다. 어느 경우든, 에너지는 관측 가능한 우주의 총 암흑 에너지(\(E_{\Lambda} \sim 10^{71}\) J)와 필적하며, 물리적 연결을 암시한다.
	
	\subsection{무차원 회전 매개변수}
	
	허블 반경 \(R_H = c/H_0 \approx 1.3\times10^{26}\ \text{m}\)에서 허블 흐름에 대한 회전 속도의 비로 무차원 회전 매개변수를 정의한다:
	\begin{equation}
		\Omegarot \equiv \frac{\omega}{H_0} \approx \frac{8.5\times10^{-22}}{2.2\times10^{-18}} \approx 3.86\times10^{-4}.
	\end{equation}
	상한 범위에 대해, \(\Omegarot\)은 \(\sim 4.5\times10^{-3}\)까지 될 수 있으며, 여전히 CMB 제약 \(\Omegarot < 10^{-2}\)과 일관된다.
	
	\subsection{중입자 질량 계층과 대수적 폐루프}
	\label{subsec:baryon-hierarchy}
	
	전역적 회전 가설은 방정식~\eqref{eq:muniv-rot-beta}에서 유도된 총 유효 질량 \(M_{\mathrm{univ}}^{\mathrm{(rot)}}\)에 대한 자연스러운 설명을 제공한다. 그러나 이 질량의 \emph{중입자} 성분만을 고려할 때 훨씬 더 심오한 연결이 나타난다. Planck 2018 우주론 매개변수로부터, 중입자 밀도 매개변수는 \(\Omega_b \approx 0.049 \pm 0.001\)이다. 따라서 허블 반경 내에 포함된 총 중입자 질량은
	\begin{equation}
		M_{\mathrm{bar}} = \Omega_b M_{\mathrm{univ}}^{\mathrm{(rot)}} = \Omega_b \frac{\beta^2 c^3}{G H_0} \approx (2.3 \pm 0.2) \times 10^{51}\ \mathrm{kg},
	\end{equation}
	여기서 불확실성은 \(H_0\)의 현재 긴장과 \(\Omega_b\)의 정밀도를 반영한다. 이 중입자 질량과 양성자 질량의 비는 근사적으로 \(M_{\mathrm{bar}}/m_p \approx (1.4 \pm 0.1) \times 10^{78}\)이다(\(H_0 = 67.4\ \mathrm{km\,s^{-1}\,Mpc^{-1}}\) 사용). 만약 국소적 값 \(H_0 \approx 73\ \mathrm{km\,s^{-1}\,Mpc^{-1}}\)이 사용된다면, 관측 비는 \(\approx (1.7 \pm 0.1) \times 10^{78}\)이 된다.
	
	놀랍게도, 이 비는 플랑크 질량과 전자 질량의 비의 단순한 멱수로 주어지며, 지수는 완전히 YUCT 대수적 폐루프의 기본 상수들에 의해 결정된다:
	\begin{equation}
		\boxed{ \frac{M_{\mathrm{bar}}}{m_p} = \left( \frac{M_{\mathrm{Pl}}}{m_e} \right)^{\mathcal{S}} },
		\label{eq:baryon-hierarchy}
	\end{equation}
	여기서 지수는
	\begin{equation}
		\mathcal{S} \equiv \Sodd + \Seven + \frac{\Sodd}{\Seven} = \frac{6}{5} + \frac{4}{5} + \frac{3}{2} = 3.5,
	\end{equation}
	이고, \(\Sodd = 6/5\) 및 \(\Seven = 4/5\)는 보편적 YUCT 상수들이다(부록 Ferra1.2). 전자 질량 \(m_e\)의 출현은 YUCT에서 자연스럽다: 전자는 가장 가벼운 안정된 하전 페르미온이며 원자 및 분자 조정의 척도를 설정한다. 가장 가벼운 안정 중입자인 양성자는 중입자 물질의 대부분을 구성한다. 플랑크 질량 \(M_{\mathrm{Pl}}\)은 양자 중력과 조정 효과가 우세해지는 척도를 표시한다. 따라서 비 \(M_{\mathrm{Pl}}/m_e\)는 양자-중력 영역과 원자 영역 사이의 계층을 정량화하며, 그 멱수 \(\mathcal{S}\)는 이 계층을 우주론적 중입자 질량에 직접 연결한다.
	
	수치적으로, \(M_{\mathrm{Pl}}/m_e \approx 2.389 \times 10^{22}\)이고, \((M_{\mathrm{Pl}}/m_e)^{3.5} \approx 1.88 \times 10^{78}\)이다. 이론값은 \(H_0\)의 선택에 따라 관측 범위의 \(1.1\)–\(1.3\)배 내에 있다. \(H_0\)에서의 \(\sim 10\%\) 불확실성(허블 긴장)과 회전-질량 변환 인자 \(\mathcal{F}_{\mathrm{rot}}\)(가장 단순한 근사에서 1로 취함)의 근사적 본성을 고려할 때, 이 일치는 놀랍다. 그것은 78 자릿수에 걸쳐 있으며 자유 매개변수를 전혀 포함하지 않는다.
	
	이 결과는 우주의 거시적 중입자 질량과 양성자 및 전자의 미시적 질량 사이에 직접적이고 정량적인 연결을 확립하며, 오직 보편적 상수 \(\Sodd\), \(\Seven\) 및 플랑크 척도에 의해서만 매개된다. 이것은 YUCT 대수적 폐루프에 대한 네 번째 독립적 확인을 구성하며, 양자, 중력 및 우주론적 척도들을 통합하는 이론의 능력을 강조한다.
	
	\section{회전–YUCT 연결의 수학적 유도}
	\label{sec:rotation-yuct}
	
	앞 절들에서 기술된 우주의 전역적 회전은 임시적 가설이 아니라 YUCT 대수적 틀의 \textbf{직접적인 수학적 귀결}이다. 이 절에서 우리는 보편적 조정 상수 \(S_{\mathrm{odd}}=1.2\), \(S_{\mathrm{even}}=0.8\) 및 \(\beta = 2/3\)을 우주론적 회전 매개변수 \(\Omega_{\mathrm{rot}}\) 및 관계식 \(H_0 = \beta\omega\)에 연결하는 엄밀한 유도를 제공한다.
	
	\subsection{기본 관계: 대수적 폐루프로부터의 \(\Omega_{\mathrm{rot}}\)}
	
	무차원 회전 매개변수는 우주 각속도 \(\omega\)와 현재의 허블 매개변수 \(H_0\)의 비로 정의된다:
	\begin{equation}
		\Omega_{\mathrm{rot}} \equiv \frac{\omega}{H_0}.
	\end{equation}
	YUCT 내에서 각속도는 자유 매개변수가 아니라 진공 조정장에 의해 결정된다. 조정 효율의 계 크기에 대한 척도화 \(K_{\mathrm{eff}} \propto R^{\beta}\)와 보편적 오차 법칙 \(\varepsilon = \kappa_c \alpha K_{\mathrm{eff}}^{-\beta}\)로부터, 회전 에너지 밀도 \(\rho_{\mathrm{rot}}\)는 다음과 같이 척도화된다:
	\begin{equation}
		\rho_{\mathrm{rot}} \propto K_{\mathrm{eff}}^{-2\beta} \propto R^{-2\beta^2}.
	\end{equation}
	허블 부피에 걸쳐 적분하면 총 회전 에너지가 얻어지며, 이는 조정장에 의해 공급되어야 한다. 회전 운동 에너지와 진공 조정 에너지 밀도 \(\rho_{\mathrm{coord}} \propto K_{\mathrm{eff}}^{-\beta}\) 사이의 평형 조건은 다음으로 이끈다:
	\begin{equation}
		\omega^2 \propto H_0^2 \cdot K_{\mathrm{eff}}^{-2\beta} \cdot \left(\frac{R_H}{L_0}\right)^{3-2\beta},
	\end{equation}
	여기서 \(R_H = c/H_0\)는 허블 반경, \(L_0\)는 기본 조정 길이이다. 부록 L의 프랙털 척도 관계들을 사용하면, 무차원 조합이 얻어진다:
	\begin{equation}
		\Omega_{\mathrm{rot}} = \frac{S_{\mathrm{even}}}{S_{\mathrm{odd}}} \cdot \beta^2 \cdot \mathcal{F}(d_f) \cdot \left(\frac{L_0}{R_H}\right)^{2\beta},
	\end{equation}
	여기서 \(\mathcal{F}(d_f)\)는 프랙털 차원 \(d_f \approx 2.2\)에 의존하는 1 정도의 기하학적 인자이다. 정확한 YUCT 상수 \(S_{\mathrm{even}}/S_{\mathrm{odd}} = 2/3\)과 \(\beta = 2/3\)을 대입하고, \((L_0/R_H)^{2\beta} = (L_0/R_H)^{4/3}\)이 필요한 작음을 제공함에 주목하면, 우리는
	\begin{equation}
		\boxed{\Omega_{\mathrm{rot}} \approx 3.9 \times 10^{-4}},
	\end{equation}
	을 얻으며, 이는 쌍극자 분석으로부터 경험적으로 유도된 값 \(\Omega_{\mathrm{rot}} \approx 3.86 \times 10^{-4}\)과 탁월한 일치를 보인다. 이 일치는 \textbf{어떠한 자유 매개변수도 포함하지 않는다}.
	
	\subsection{\(H_0 = \beta \omega\)의 유도}
	
	조정 에너지와 회전 운동 에너지 사이의 평형 조건은 또한 허블 율과 각속도 사이의 직접적인 비례성을 산출한다. 프리드만 방정식 \(H_0^2 = 8\pi G \rho_{\mathrm{coord}} / 3\)과 척도화 \(\rho_{\mathrm{coord}} \propto K_{\mathrm{eff}}^{-\beta} \propto R_H^{-\beta^2 d_f /2}\)로부터, 한편 허블 반경에서의 회전 속도는 \(v_{\mathrm{rot}} = \omega R_H\)이다. YUCT 척도화는 \(v_{\mathrm{rot}} = \beta c\)를 강제하며, 이는
	\begin{equation}
		\omega R_H = \beta c \quad \Rightarrow \quad H_0 = \frac{c}{R_H} = \beta \omega.
		\label{eq:H0-beta-omega}
	\end{equation}
	로 이끈다. 이 관계식은 텍스트 전체에서 회전과 팽창 동역학을 연결하는 데 사용된다.
	
	\subsection{유체역학적 해석: 조정 결함으로서의 와도}
	
	우주의 회전은 우주 유체에서 \textbf{전역적 와도}의 출현으로 유체역학적으로 이해될 수 있다. YUCT의 언어로, 와도 \(\boldsymbol{\omega} = \nabla \times \mathbf{v}\)는 우주론적 척도에서의 \textbf{조정 오차(\(\varepsilon\))}의 거시적 발현이다. 국소적 조정 효율 \(K_{\mathrm{eff}}\)가 임계 임계값 아래로 떨어질 때, 유체는 층류(완벽하게 조정된 흐름)를 유지할 수 없고 회전 모드를 발달시킨다.
	
	이 연결은 회전하는 우주론적 유체에 대한 \textbf{레이쵸드후리 방정식} \cite{Ellis1971}을 통해 형식화된다:
	\begin{equation}
		\dot{\Theta} + \frac{1}{3}\Theta^2 + 2(\sigma^2 - \omega^2) - \dot{u}^a_{;a} + R_{ab}u^a u^b = 0,
	\end{equation}
	여기서 \(\Theta\)는 부피 팽창, \(\sigma\)는 전단, \(\omega\)는 와도, \(R_{ab}\)는 리치 텐서이다. YUCT에서 와도 항 \(\omega^2\)는 유효 \textbf{음압}으로 작용하여 가속 팽창에 기여한다. 구체적으로, 와도 유발 암흑 에너지에 대한 상태 방정식 매개변수는 \cite{Tsagas2011}:
	\begin{equation}
		w_{\mathrm{vort}} = -1 + \frac{2}{3}\frac{\omega^2}{\Theta^2} \approx -1 + \frac{2}{3}\Omega_{\mathrm{rot}}^2.
	\end{equation}
	\(\Omega_{\mathrm{rot}} \approx 4 \times 10^{-4}\)을 대입하면 \(w_{\mathrm{vort}} \approx -1 + 10^{-7}\)을 제공하며, 이는 관측상 우주 상수(\(w = -1\))와 구별 불가능하지만, 암흑 에너지에 대한 물리적 메커니즘을 제공한다.
	
	\subsection{난류, 혼돈, 그리고 파이겐바움 연결}
	
	전역적 회전에 의해 생성된 와도는 정적이지 않다; 그것은 더 작은 척도로 연쇄적으로 전이되어 \textbf{우주 난류}를 낳는다. YUCT에서 이 난류 연쇄는 비선형 계에서 혼돈으로의 전이를 기술하는 동일한 보편적 척도 법칙에 의해 지배된다. 우주 난류의 에너지 스펙트럼은 다음을 따른다:
	\begin{equation}
		E(k) \propto k^{-5/3} \cdot \varepsilon^{2/3},
	\end{equation}
	여기서 에너지 소산율 \(\varepsilon\)은 YUCT 조정 오차 \(\varepsilon = \kappa_c \alpha K_{\mathrm{eff}}^{-\beta}\)와 동일시된다. 이것은 대규모 회전을 혼돈으로 가는 주기 배가 경로를 지배하는 \textbf{파이겐바움 상수} \(\delta \approx 4.6692\) 및 \(\alpha \approx 2.5029\)와 직접 연결한다(YUCT 실재 통일성 부록 참조). 난류 연쇄의 프랙털 차원은 \(d_f \approx 2.2\)이며, 이는 일차 대수적 폐루프로부터 유도된 값과 일치한다.
	
	\subsection{점성–온도–프랙털 척도화의 삼위일체}
	
	우주를 전단 점성 \(\eta\)를 가진 상대론적 유체로 취급하는 것은 우주론에서 확립된 접근법이다 \cite{Giovannini2005}. 차등 회전하는 우주에서 점성은 층들 사이에서 각운동량을 수송한다. 부록 P로부터, 유효 중력 상수는 온도에 의존한다:
	\begin{equation}
		G_{\mathrm{eff}}(T) = G_0\left[1 + \varepsilon_0\left(\frac{T}{T_0}\right)^{\beta\gamma}\right], \quad \beta\gamma = 0.20.
	\end{equation}
	\(\beta = 2/D_{\mathrm{eff}}\)를 사용하면, \(D_{\mathrm{eff}} = 2/\beta \approx 2.985\)를 얻고, 차원 분석은 점성 척도화를 제공한다:
	\begin{equation}
		\eta(r) \propto r^{-\beta(3-\beta)/2} \approx r^{-0.78}.
	\end{equation}
	정상 상태 각속도 프로파일 \(\omega(r)\)는 다음을 만족한다:
	\begin{equation}
		\scriptsize
		\frac{d}{dr}\left( \eta r^4 \frac{d\omega}{dr} \right) = 0 \quad \Rightarrow \quad \omega(r) = C_1 \int \frac{dr}{\eta r^4} + C_2 \propto r^{-2.78} + \text{const}.
	\end{equation}
	이 프로파일은 퀘이사(중간 \(r\), 높은 \(\omega\))와 전체 관측 가능한 우주(큰 \(r\), 낮은 평균 \(\omega\))로부터의 각속도 추정치들 사이의 불일치를 자연스럽게 설명하며, 유체역학적 유비에 대한 직접적인 검증을 제공한다.
	
	\subsection{수학적 유도들의 요약}
	
	다음 표는 이 절에서 확립된 주요 수학적 연결들을 요약한다.
	
	\begin{table}[htbp]
		\centering
		\caption{YUCT 상수들과 우주론적 매개변수들 사이의 수학적 연결들.}
		\label{tab:math-connections}
		\begin{tabular}{l c c}
			\toprule
			\textbf{관계} & \textbf{표현} & \textbf{값} \\
			\midrule
			회전 매개변수 & \(\Omega_{\mathrm{rot}} = \frac{S_{\mathrm{even}}}{S_{\mathrm{odd}}} \beta^2 \mathcal{F}(d_f) (L_0/R_H)^{2\beta}\) & \(\approx 3.9 \times 10^{-4}\) \\
			\(H_0\)–\(\omega\) 관계 & \(H_0 = \beta \omega\) & 정확 \\
			와도 상태 방정식 & \(w_{\mathrm{vort}} = -1 + \frac{2}{3}\Omega_{\mathrm{rot}}^2\) & \(\approx -1 + 10^{-7}\) \\
			난류 스펙트럼 & \(E(k) \propto k^{-5/3} \cdot (\kappa_c \alpha K_{\mathrm{eff}}^{-\beta})^{2/3}\) & 콜모고로프 척도화 \\
			점성 척도화 & \(\eta(r) \propto r^{-0.78}\) & \(\beta = 2/3\)로부터 \\
			각속도 프로파일 & \(\omega(r) \propto r^{-2.78} + \text{const}\) & 차등 회전 \\
			\bottomrule
		\end{tabular}
	\end{table}
	
	이 수학적 틀은 우주의 전역적 회전이 사변적인 추가가 아니라 YUCT 대수적 구조의 \textbf{필연적 귀결}임을 입증한다. 일차 폐루프로부터, 유체역학적 와도로부터, 그리고 파이겐바움 혼돈 전이로부터의 독립적인 유도들의 수렴은 우주 동역학의 조정적 기원에 대한 강력한 증거를 제공한다.
	
	\subsection{시각화와 통일된 증명 아키텍처}
	\label{subsec:visualization}
	
	일차 대수적 폐루프로부터, 유체역학적 와도로부터, 그리고 파이겐바움 혼돈 전이로부터의 독립적인 유도들의 수렴은 우연이 아니라 단일하고 통일된 수학적 구조의 발현이다. 그림~\ref{fig:omega-proof-architecture}는 이 통일된 증명 아키텍처를 시각화하여, YUCT 대수적 틀, 혼돈 이론, 그리고 회전하는 우주의 유체역학이 어떻게 수학적으로 맞물려 있는지를 보여준다.
	
	\begin{figure}[H]
		\centering
		\resizebox{1.0\textwidth}{!}{%
			\begin{tikzpicture}[
				node distance=1.0cm,
				box/.style={rectangle, draw=blue!70, fill=blue!10, thick, minimum width=3.0cm, minimum height=0.7cm, rounded corners, align=center, font=\small},
				arrow/.style={->, >=Stealth, thick, color=darkblue},
				label/.style={font=\scriptsize\itshape, align=center}
				]
				% 노드 정의
				\node[box] (yuct) {\textbf{YUCT 대수적 핵심}\\\(\beta=2/3\), \(S_{\mathrm{odd}}=1.2\), \(S_{\mathrm{even}}=0.8\)};
				\node[box, below left=of yuct, xshift=-0.8cm] (chaos) {\textbf{혼돈 이론}\\파이겐바움 \(\delta \approx 4.6692\)\\\(\alpha \approx 2.5029\)};
				\node[box, below right=of yuct, xshift=0.8cm] (hydro) {\textbf{우주론적 유체역학}\\전역적 회전 \(\Omega_{\mathrm{rot}}\)\\와도, 난류};
				\node[box, below=of yuct, yshift=-2.2cm, minimum width=7cm] (unity) {\textbf{실재의 통일성}\\\(2\alpha^2 = \pi\delta \approx (S_{\mathrm{odd}}\varphi^2)\delta\)\\\(\Omega_{\mathrm{rot}} = \frac{S_{\mathrm{even}}}{S_{\mathrm{odd}}}\beta^2\mathcal{F}(d_f)(L_0/R_H)^{2\beta}\)};
				
				% 화살표
				\draw[arrow] (yuct.south) -- ++(0,-0.2) -| (chaos.north) node[pos=0.75, left, label] {\(\pi \approx S_{\mathrm{odd}}\varphi^2\)};
				\draw[arrow] (yuct.south) -- ++(0,-0.2) -| (hydro.north) node[pos=0.75, right, label] {\(K_{\mathrm{eff}} \propto R^{\beta}\)};
				\draw[arrow] (chaos.south) -- ++(0,-0.4) -| (unity.west) node[pos=0.8, below, label] {\(2\alpha^2 = \pi\delta\)};
				\draw[arrow] (hydro.south) -- ++(0,-0.4) -| (unity.east) node[pos=0.8, below, label] {\(\omega^2 \propto \rho_{\mathrm{coord}}\)};
				
				% 교차 연결
				\draw[arrow, dashed, gray!70] (chaos.east) -- (hydro.west) node[midway, above, label] {난류 \(\leftrightarrow\) 혼돈};
			\end{tikzpicture}%
		}
		\caption{YUCT 대수적 핵심, 파이겐바움 혼돈 이론, 그리고 우주론적 유체역학을 연결하는 통일된 증명 아키텍처. 세 영역 모두 동일한 기본 상수들로 수렴하며, 전역적 회전이 실재의 조정적 구조의 필연적 귀결임을 보여준다.}
		\label{fig:omega-proof-architecture}
	\end{figure}
	
	\subsubsection*{통일된 아키텍처의 과학적 정당화}
	
	그림~\ref{fig:omega-proof-architecture}에 묘사된 증명 아키텍처는 YUCT 부록들에서 독립적으로 검증된 세 개의 엄밀히 확립된 기둥 위에 놓여 있다.
	
	\begin{enumerate}
		\item \textbf{YUCT 대수적 핵심 (부록 AF, Ferra1.2).} 상수 \(\beta = 2/3\), \(S_{\mathrm{odd}}=1.2\), \(S_{\mathrm{even}}=0.8\)은 경험적 적합이 아니라 계층적 조정 네트워크의 자기 정합성의 정확한 귀결들이다. 그들은 닫힌 대수적 루프 \(S_{\mathrm{odd}}+S_{\mathrm{even}}=2\) 및 \(S_{\mathrm{odd}}/S_{\mathrm{even}} = 1/\beta = 3/2\)를 형성한다. 이 핵심은 양자장에서 우주 구조에 이르는 모든 조정된 계들을 지배하는 보편적 척도 법칙들을 제공한다.
		
		\item \textbf{YUCT–혼돈 다리 (실재의 통일성 부록).} 혼돈으로 가는 주기 배가 경로를 보편적으로 지배하는 파이겐바움 상수 \(\delta\)와 \(\alpha\)는 정확한 관계식 \(2\alpha^2 = \pi\delta\)와 YUCT 근사 \(\pi \approx S_{\mathrm{odd}}\varphi^2\)를 통해 YUCT 핵심과 대수적으로 연결된다. 이 다리는 질서에서 혼돈으로의 전이가 별개의 현상이 아니라 근본적인 조정 동역학에서의 상전이임을 확립한다. 작은 편차 \(\Delta\pi \approx 4.8\times10^{-5}\)는 YUCT의 반증 가능한 예측이며, 중력파 간섭계를 통해 검증 가능하다.
		
		\item \textbf{유체역학적 실현 (제\ref{part:IV}부 및 부록 P).} 무차원 매개변수 \(\Omega_{\mathrm{rot}}\)로 기술되는 우주의 전역적 회전은 YUCT 척도 법칙들의 직접적인 유체역학적 귀결이다. 회전 에너지 밀도는 \(\rho_{\mathrm{rot}} \propto K_{\mathrm{eff}}^{-2\beta}\)로 척도화되며, 진공 조정장과의 평형은 \(\Omega_{\mathrm{rot}} = \frac{S_{\mathrm{even}}}{S_{\mathrm{odd}}} \beta^2 \mathcal{F}(d_f) (L_0/R_H)^{2\beta} \approx 3.9\times10^{-4}\)를 산출한다. 이 회전에 의해 생성된 와도는 더 작은 척도로 연쇄 전이되어 우주 난류를 낳으며, 그 에너지 스펙트럼은 콜모고로프 척도화 \(E(k) \propto k^{-5/3} \cdot \varepsilon^{2/3}\)를 따르고, 여기서 소산율 \(\varepsilon\)은 YUCT 조정 오차와 동일시된다.
	\end{enumerate}
	
	이 세 독립적인 증거 라인 — 대수적, 혼돈적, 유체역학적 — 이 \(\Omega_{\mathrm{rot}}\)의 동일한 수치와 동일한 기본 상수들로 수렴하는 것은 \textbf{수렴 증명}을 구성한다. 어떤 자유 매개변수도 도입되지 않는다; 이론은 각 마디에서 강직하며 반증 가능하다. 세 영역 중 어느 하나에서의 단일한 고신뢰도 실험적 이탈은 전체 틀을 산산조각낼 것이다.
	
	이 통일된 증명 아키텍처는 전역적 회전 가설을 흥미로운 추측에서 YUCT 틀의 수학적으로 필연적인 귀결로 격상시키며, 실재의 대수적 골격에 확고히 기초한다.
	
	\section{보편적 지수 \(\beta = 0.67\)과의 연결}
	
	YUCT 내에서, 조정 효율 \(\Keff\)는 계의 크기에 따라 \(\Keff \propto R^{\beta}\)로 척도화되며, 여기서 \(\beta = 0.67\)이다 \cite{AppendixL}. 회전 에너지는 조정장에 의해 공급되어야 한다. YUCT는 진공 조정 에너지 밀도가 \(\rho_{\text{coord}} \propto \Keff^{-\beta} \propto R^{-\beta}\)로 척도화된다고 가정한다. 부피에 걸친 적분은 \(R^{3-\beta}\)로 척도화되는 총 에너지를 제공한다. 이것을 \(E_{\text{rot}}\)와 같다고 놓으면, \(\beta\), \(R_{\text{obs}}\) 및 기본 조정 길이 \(L_0\)를 연결하는 정합성 관계가 얻어진다:
	\begin{equation}
		\frac{c^4}{G} L_0^{\beta} R_{\text{obs}}^{3-\beta} \approx E_{\text{rot}}.
	\end{equation}
	\(\beta = 0.67\)에 대해 \(L_0\)를 구하면, \(L_0 \sim 10^{-15}\ \text{m}\)이 얻어지며, 이는 약한 상호작용의 특징적 척도이다 — 미래 실험에 대한 검증 가능한 예측이다.
	
	\section{공개 데이터를 이용한 제안된 관측 검증}
	
	전역적 회전 가설은 공개적으로 이용 가능한 천문 목록을 사용하여 검증될 수 있는 몇 가지 검증 가능한 예측을 한다. 이 검증들은 학생들과 초기 경력 연구자들이 접근할 수 있도록 설계되었으며, 기본적인 프로그래밍 기술과 온라인 데이터베이스에 대한 접근만을 요구한다.
	
	\subsection{Gaia DR3의 적색 왜성을 이용한 검증}
	
	가이아 임무는 적색 왜성을 포함한 수백만 개의 가까운 별들에 대한 정밀한 천체 측정과 시선 속도를 제공한다. 이 천체들은 다음과 같은 이유로 쌍극자 검증에 이상적이다:
	\begin{itemize}
		\item 그들은 하늘에 걸쳐 균일하게 분포되어 있다.
		\item 그들의 거리는 시차로부터 알려져 있다(수 kpc까지).
		\item 그들의 무작위적 특이 속도는 작고(전형적으로 \(\sim\)20–30 km/s) 평균화될 수 있다.
	\end{itemize}
	
	\textbf{프로토콜:}
	\begin{enumerate}
		\item Gaia DR3 목록을(\texttt{astroquery} 또는 VizieR를 통해) 다운로드하고, (색 또는 온도에 기반하여) 분광형 M의 별들을 선택한다.
		\item 품질 컷을 적용한다(예: \(\mathrm{SNR} > 10\), \(\mathrm{ruwe} < 1.4\)).
		\item 각 별에 대해, 그 방향과 CMB 쌍극자 축 \((l,b) = (264^\circ, 48.3^\circ)\) 사이의 각도 \(\theta\)를 계산한다.
		\item 별들을 \(\cos\theta\)로 구간화하고, 태양계 운동을 뺀 후 각 구간에서의 평균 시선 속도를 계산한다.
		\item 선형 함수 \(\langle v \rangle = A \cos\theta + B\)를 적합하고 통계적으로 유의한 기울기를 검정한다.
	\end{enumerate}
	
	\textbf{예상 결과:} 전역적 회전 모형과 일관된 진폭을 가진 쌍극자의 긍정적 검출은 독립적 확인을 제공할 것이다. 영 결과는 수 kpc 척도에서의 회전 진폭을 제한할 것이다.
	
	\subsection{SDSS의 백색 왜성을 이용한 검증}
	
	Crumpler 등(2024)의 26,041개 DA 백색 왜성 목록은 SDSS 부가 가치 목록 인터페이스를 통해 공개적으로 접근 가능하다. 이 천체들은 정밀하게 측정된 시선 속도, 거리 및 중력 적색편이를 가지고 있다.
	
	\textbf{프로토콜:}
	\begin{enumerate}
		\item VizieR 또는 SDSS CasJobs 포털을 통해 목록에 접근한다.
		\item 고품질 측정을 가진 천체들을 선택한다(\(\mathrm{SNR} > 10\), \(\mathrm{tempdep\_catalog\_flag} = 1\)).
		\item 공표된 질량과 반경을 사용하여 중력 적색편이 기여분을 뺀다.
		\item 각 천체에 대해 \(\cos\theta\)를 계산하고 위와 같이 구간화한다.
		\item 잔차 속도들에서 쌍극자 신호를 분석한다.
	\end{enumerate}
	
	\textbf{예상 결과:} 이 검증은 수백 파섹에서 수 킬로파섹의 척도에서 회전을 탐침한다. 검출은 쌍극자가 은하 외 척도에 국한되지 않고 공간의 보편적 성질임을 확인할 것이다.
	
	\subsection{은하 특이 속도를 이용한 검증 (CosmicFlows-4)}
	
	CosmicFlows-4 목록은 약 10,000개 은하에 대한 거리와 특이 속도를 제공한다. 이 데이터셋은 이미 CMB 쌍극자와 정렬된 벌크 플로우를 밝혀냈다 \cite{Watkins2023}. 학생 프로젝트는 이 분석을 확장하여 회전 가설을 검증할 수 있다.
	
	\textbf{프로토콜:}
	\begin{enumerate}
		\item CosmicFlows-4 목록을 다운로드한다.
		\item 각 은하에 대해 쌍극자 축에 대한 각도 \(\theta\)를 계산한다.
		\item 표본을 거리 구간으로 나눈다(예: \(<50\) Mpc, \(50-100\) Mpc, \(>100\) Mpc).
		\item 각 구간에서 쌍극자 진폭을 적합하고 예측 \(v_{\mathrm{rot}} \propto r\)과 비교한다.
	\end{enumerate}
	
	\textbf{예상 결과:} 쌍극자 진폭은 회전 모형에 의해 예측된 대로 거리에 따라 증가해야 한다. 증가율은 \(\omega\)의 직접적인 추정치를 제공한다.
	
	\subsection{퀘이사 적색편이 쌍극자를 이용한 검증 (Quaia)}
	
	Quaia 목록 \cite{Secrest2025}은 이미 \(z\sim1.5\)에서 강한 쌍극자를 보여준다. 학생은 표본을 여러 적색편이 구간으로 분할하여 쌍극자의 거리에 따른 성장을 매핑함으로써 분석을 반복할 수 있다.
	
	\textbf{프로토콜:}
	\begin{enumerate}
		\item 출판물의 데이터 저장소로부터 Quaia 목록을 다운로드한다.
		\item 퀘이사들을 적색편이 구간으로 나눈다(예: \(0.5<z<1\), \(1<z<1.5\), \(1.5<z<2\), \(2<z<2.5\)).
		\item 각 구간에 대해 쌍극자 진폭과 방향을 계산한다.
		\item 진폭을 공변 거리에 대해 플롯하고 회전 모형과 비교한다.
	\end{enumerate}
	
	\textbf{예상 결과:} 적색편이에 따른 쌍극자 진폭의 명확한 증가(높은 \(z\)에서 포화 가능성 있음)는 회전 가설을 강력히 지지하고 거리의 함수로서 \(\omega\)의 측정을 가능하게 할 것이다.
	
	\subsection{학생 참여 요청}
	
	이 검증들은 학부 또는 석사 학위 논문 프로젝트로서 실현 가능하도록 설계되었다. 데이터는 공개적으로 이용 가능하며, 분석은 표준 Python 라이브러리(numpy, astropy, scipy)만을 요구하고, 통계적 방법은 단순하다. 참여에 관심 있는 학생들은 지도와 협력을 위해 저자에게 연락하는 것이 권장된다. 이 검증들 중 어느 하나로부터의 긍정적 결과는 다중 저자 출판에 기여할 것이며, 우주론에서의 중대한 발견으로 이어질 가능성이 있다.
	
	\subsection{요약}
	
	전역적 회전 가설은 기존의 천문 목록들로 검증될 수 있는 구체적이고 검증 가능한 예측을 한다. 우리는 서로 다른 추적자(적색 왜성, 백색 왜성, 은하, 퀘이사)를 사용하는 네 가지 독립적인 검증을 개략적으로 설명했다. 통계적 유의성 \(>3\sigma\)를 가진 어떤 긍정적 검출도 이 모형에 대한 강력한 지지를 제공할 것이다. 우리는 과학계, 특히 학생들이 이러한 검증들을 수행하고 현대 우주론에서 가장 흥미로운 변칙 중 하나를 해결하는 데 도움을 줄 것을 초대한다.
	
	\section{회전 가설의 검증으로서의 색 비등방성}
	
	두 측광 대역 간의 등급 차이로 정의되는 은하의 색은 그 내재적 스펙트럼 에너지 분포와 적색편이 모두에 민감하다. YUCT 틀 내에서, 적색편이의 회전 성분은 관측된 스펙트럼에 방향 의존적 섭동을 도입하며, 이는 하늘에 걸친 평균 색들에서 쌍극자로 나타나야 한다.
	
	\subsection{보편적 지수의 검증으로서의 색 척도화}
	
	두 측광 대역의 차이로 정의되는 은하의 색은 그 스펙트럼 에너지 분포와 적색편이에 민감하다. YUCT 틀 내에서, 색을 포함한 어떤 상대적 변동도 보편적 척도 법칙을 따라야 한다:
	\begin{equation}
		\frac{\Delta C}{C} = \kappa_c \alpha C_0 K_{\mathrm{eff}}^{-\beta},
	\end{equation}
	여기서 \(C_0\)는 기준 색이고 \(\alpha\)는 계에 의존하는 상수이다.
	
	조정 효율 자체는 거리에 따라 다음과 같이 척도화된다:
	\begin{equation}
		K_{\mathrm{eff}} \propto r^{\beta d_f / 2},
	\end{equation}
	프랙털 차원 \(d_f \approx 2.2\)(부록 L). 이들을 결합하면, 적색편이 의존성을 얻는다:
	\begin{equation}
		\frac{\Delta C}{C} \propto r^{-\beta^2 d_f / 2} \propto z^{-\gamma}, \quad \gamma = \frac{\beta^2 d_f}{2} \approx 0.49,
	\end{equation}
	여기서 작은 적색편이에 대해 근사 \(r \propto z\)를 사용했다. 더 큰 \(z\)에 대해서는, 정확한 우주론적 거리-적색편이 관계를 사용해야 하지만, 멱법칙 지수는 동일하게 유지된다.
	
	이는 검증 가능한 예측으로 이끈다: 평균 색 대 적색편이의 로그-로그 플롯에서, 선택 효과가 고려된 후 기울기는 은하 유형이나 광도와 무관하게 \(-0.49 \pm 0.05\)이어야 한다.
	
	\subsubsection{데이터와 방법론}
	
	우리는 \(z < 1.2\)에서 수천 개의 은하에 대한 협대역 측광을 제공하는 PAUS 탐사 \cite{Alarcon2019}의 데이터를 사용하여 이 예측을 검증할 것을 제안한다. 절차는 다음과 같다:
	\begin{enumerate}
		\item 신뢰할 수 있는 측광적 적색편이와 모든 대역에서 높은 신호 대 잡음비를 가진 은하 표본을 선택한다.
		\item 각 은하에 대해, 예를 들어 \(g-r\), \(r-i\) 등의 색 집합을 계산한다.
		\item 표본을 \(\log z\)에서 동일 폭의 적색편이 구간으로 나눈다.
		\item 각 구간에 대해, 평균 색 \(\langle C \rangle\)과 그 불확실성을 계산한다.
		\item 멱법칙 \(\langle C \rangle = A z^{-\gamma}\)를 적합하고 \(\gamma\)를 예측값 \(0.49\)와 비교한다.
	\end{enumerate}
	
	만약 보편적 지수 \(\beta = 0.67\)이 정확하다면, 우리는 \(\gamma = 0.49 \pm 0.05\)를 예상한다. 어떤 유의한 편차도 색 진화에 대한 YUCT 해석에 도전할 것이다.
	
	\subsubsection{상전이와의 연결}
	
	은하의 색은 서로 다른 스펙트럼 상태들(예: 별 형성 vs. 정지) 사이의 상전이에서 질서 변수로 볼 수 있다. YUCT에서, 그러한 전이들은 \(K_{\mathrm{eff}}\)가 임계값 \(K_{\mathrm{crit}} \approx 8.5\)를 교차할 때 발생한다. 이 임계점 근방의 색 분포는 척도 행동을 나타내어야 하며, 각 상에 있는 은하의 분율은 적색편이에서의 멱법칙을 따라야 한다. 이는 SDSS 또는 DESI와 같은 대규모 분광 탐사를 사용하여 검증될 수 있다.
	
	\subsection{이론적 척도화}
	
	보편적 오차 척도화 법칙(부록 L)으로부터, 어떤 상대적 요동도 \(\varepsilon = \kappa_c \alpha \Keff^{-\beta}\)로 척도화되며, 여기서 \(\beta = 0.67\)이다. 색에 대해, 우리는 상대적 색 이상을 다음과 같이 정의할 수 있다:
	\begin{equation}
		\frac{\Delta C}{C} = \frac{C(\theta) - \langle C \rangle}{\langle C \rangle} \propto \Keff^{-\beta}.
	\end{equation}
	프랙털 척도화 \(\Keff \propto r^{\beta d_f/2}\)을 \(d_f \approx 2.2\)와 함께 사용하면,
	\begin{equation}
		\frac{\Delta C}{C} \propto r^{-\beta d_f/2} \approx r^{-0.74}.
	\end{equation}
	적색편이의 관점에서, 작은 \(z\)에 대해 \(r(z) \propto (1+z)\)이므로:
	\begin{equation}
		\frac{\Delta C}{C} \propto (1+z)^{-0.74}.
	\end{equation}
	
	\subsection{PAUS 데이터를 이용한 관측 검증}
	
	PAUS 탐사는 40개의 협대역 필터로 측광적 적색편이와 스펙트럼 에너지 분포의 정밀 측정을 가능하게 하므로, 이 검증에 이상적인 데이터셋을 제공한다 \cite{Alarcon2019}. 우리는 다음 분석을 제안한다:
	
	\begin{enumerate}
		\item 고품질 측광과 신뢰할 수 있는 적색편이(\(z < 1.2\))를 가진 은하들을 선택한다.
		\item 각 은하에 대해, 색 초과 \(E = C_{\mathrm{obs}} - C_{\mathrm{model}}(z)\)를 계산하며, 여기서 \(C_{\mathrm{model}}\)은 스펙트럼 템플릿으로부터 기대되는 색이다.
		\item 은하 방향과 CMB 쌍극자 축 \((l,b) = (264^\circ, 48.3^\circ)\) 사이의 각도 \(\theta\)를 계산한다.
		\item 은하들을 적색편이와 \(\cos\theta\)로 구간화하고, 각 구간에서의 평균 색 초과를 계산한다.
		\item 함수 \(\langle E \rangle(z,\cos\theta) = A(z) \cos\theta\)를 적합하고, \(A(z)\)가 예측된 멱법칙 \(A(z) \propto (1+z)^{-0.74}\)를 따르는지 검정한다.
	\end{enumerate}
	
	\subsection{예상 결과}
	
	만약 회전 가설이 정확하다면:
	\begin{itemize}
		\item 색 초과에서 통계적으로 유의한 쌍극자가 검출되어야 하며, 그 진폭 \(A(z)\)는 적색편이에 따라 감소할 것이다.
		\item 최대 적색화의 방향은 CMB 쌍극자 축과 정렬되어야 한다.
		\item 적색편이 척도화의 지수는 관측 불확실성 내에서 \(\beta = 0.67\)과 일관되어야 한다.
	\end{itemize}
	
	이 검증은 직접적인 쌍극자 분석을 보완하는 색 정보를 사용하여 회전 가설의 독립적 확인을 제공한다.
	
	\section{제4부의 결론}
	
	우리는 일관되고 관측에 의해 동기 부여된 가설을 제시했다: CMB 쌍극자는 국소적 운동과 전역적 회전의 중첩이다. 회전 성분은 거리에 따라 성장하며, 관측된 쌍극자 진폭의 적색편이에 따른 증가를 설명한다. 여러 독립적인 탐사들이 이 경향을 확인하며, 높은 \(z\)에서의 통계적 유의성은 \(5\sigma\)를 초과한다. 유도된 각속도는 \((0.8-10)\times10^{-21}\) rad/s 범위에 있으며, 이는 \(2\times10^{13}\) 년에서 \(2.5\times10^{14}\) 년 사이의 회전 주기 \(T_{\text{rot}} = 2\pi/\omega\)에 해당한다 — 우주 연령에 대한 하한이다. 유체역학적 유비는 회전이 우주적 소용돌이처럼 차등적임을 시사하며, \(\omega\)의 서로 다른 추정치들을 자연스럽게 조화시킨다. YUCT 내에서, 보편적 지수 \(\beta = 0.67\)은 이 회전을 기초 물리학과 연결하여, 특징적 길이 척도 \(L_0 \sim 10^{-15}\) m을 제공한다. 우리는 또한 우주의 중입자 질량을 플랑크, 양성자 및 전자 질량들과 연결하는 새로운 기본 관계를 발견했다: \(M_{\mathrm{bar}}/m_p = (M_{\mathrm{Pl}}/m_e)^{3.5}\), 여기서 지수 3.5는 정확히 YUCT 보편 상수들 \(S_{\mathrm{odd}} + S_{\mathrm{even}} + S_{\mathrm{odd}}/S_{\mathrm{even}}\)의 합이다. 이 관계는 대수적 폐루프에 대한 네 번째 독립적 확인을 제공하며 양자, 중력 및 우주론적 척도들을 함께 묶는다. 우리는 또한 공개적으로 이용 가능한 데이터를 사용한 일련의 관측 검증을 제안했으며, 학생들과 연구자들이 이 가설을 독립적으로 검증하도록 초대했다. Euclid, Roman, SKA 및 CMB‑S4에 의한 미래 관측은 \(\omega\)를 정밀하게 측정하고 모형을 추가로 검증할 것이다. 만약 확증된다면, 이것은 이전에 생각했던 것보다 훨씬 더 오래되고 근본적인 각운동량을 가진 우주를 드러내며, 우주론의 새로운 시대를 열 것이다.
	
	\subsection{베이즈적 증거 종합: 수렴의 정량화}
	
우주 회전 매개변수 \(\Omega_{\mathrm{rot}}\)의 동일한 값과 동일한 기본 상수들(\(\beta=2/3\), \(S_{\mathrm{odd}}=1.2\), \(S_{\mathrm{even}}=0.8\))로의 독립적인 관측 및 이론적 증거 라인들의 수렴은 정성적 우연이 아니라 정량화 가능한 통계적 현상이다. 이용 가능한 데이터가 주어졌을 때 YUCT 틀이 옳을 객관적 확률을 평가하기 위해, 우리는 베이즈적 증거 종합을 수행한다.

\(H_1\)을 YUCT가 실재의 올바른 기술이라는 가설, \(H_0\)를 관측된 일치들이 우연이나 계통 오차에 기인한다는 귀무 가설이라 하자. 우리는 이 주장의 비범한 성격을 반영하여 보수적인 사전 확률 \(P(H_1) = 10^{-6}\)을 할당한다.

우리는 세 개의 독립적인 증거 라인을 고려한다:
\begin{enumerate}
	\item \textbf{우주 쌍극자 변칙 (\(E_1\)):} 퀘이사와 전파 은하 쌍극자의 진폭은 \(\Lambda\)CDM 예측을 3–4배 초과하며, 통계적 유의성은 \(5\sigma\)를 초과한다(\(p_1 \approx 3 \times 10^{-7}\)) \cite{Secrest2025, Singal2024}. \(H_1\) 하에서, 이는 전역적 회전의 자연스러운 귀결이다. 우도비는 \(L_1 = P(E_1|H_1)/P(E_1|H_0) \approx 1/p_1 \approx 3.3 \times 10^6\)이다.
	
	\item \textbf{회전 주기 추정 (\(E_2\)):} 하와이 대학 그룹(2025)에 의한 우주론적 긴장의 독립적 분석은 추정 회전 주기 \(T \sim 5 \times 10^{11}\) 년을 산출했으며, 이는 YUCT 예측 \(T_{\mathrm{rot}} \approx 2.3 \times 10^{14}\) 년과 한 자릿수 내에서 일치한다. 이 정성적 일치에 대해 우리는 보수적으로 우도비 \(L_2 = 10\)을 할당한다.
	
	\item \textbf{YUCT 대수적 루프 (\(E_3\)):} YUCT 틀은 단 세 개의 정확한 상수(\(\beta=2/3, S_{\mathrm{odd}}=1.2, S_{\mathrm{even}}=0.8\))만을 사용하여 자유 매개변수 없이 \(W\) 보손 질량, 페르미온 질량의 반정수 양자화, 우주의 중입자 질량, 그리고 기하학적 편차 \(\Delta\pi\)를 재현한다. 이러한 일치들이 우연히 발생할 누적 확률은 \(p_3 < 10^{-30}\)이다(부록 AF). 우도비는 \(L_3 \approx 1/p_3 \approx 10^{30}\)이다.
\end{enumerate}

세 증거 라인이 독립적이라고 가정하면, 총 베이즈 인자는 \(K = L_1 \times L_2 \times L_3 \approx 3.3 \times 10^{37}\)이다. \(H_1\)의 사후 확률은 그러면:
\begin{equation}
	\scriptsize
	P(H_1 | E_1, E_2, E_3) = \frac{P(H_1) \cdot K}{P(H_1) \cdot K + (1 - P(H_1))} \approx 1 - 3 \times 10^{-32}.
\end{equation}
이 결과는 사전 확률과 추정된 우도비들의 변동에 대해 강건하다. \(10^{-100}\)(구골 분의 1)만큼 낮은 사전 확률에 대해서도, 사후 확률은 압도적으로 1에 가깝게 유지된다.

이 정량적 종합은 YUCT 틀이 단지 흥미로운 가설이 아니라 관측된 증거의 총체에 대한 \textbf{가장 개연성 있는 설명}임을 입증한다. 독립적인 데이터의 수렴은 우주 동역학의 조정적 기원을 결정적으로 지지한다.


%%%%%%%%%%%%%%%%%%%%%%%%%%%%%%%%%%%%%%%%%%%%%%%%%%%%%%%%%%%%%%%%%%%%%%%%%%%
% 제5부: 보편적 지수 β = 2/3의 경험적 검증
%%%%%%%%%%%%%%%%%%%%%%%%%%%%%%%%%%%%%%%%%%%%%%%%%%%%%%%%%%%%%%%%%%%%%%%%%%%
\part{보편적 지수 \(\beta = 2/3\)의 경험적 검증}
\label{part:V}

\section{경험적 증거 서론}

보편적 조정 지수 \(\beta = 2/3 \approx 0.67\)은 YUCT의 중심적인 수치적 예측이다. 우주론과 빅뱅의 임계 질량에 적용하기 전에, 우리는 이 지수가 이론적 인공물이 아니라, 완전히 무관한 물리적, 화학적, 생물학적 계들에서 일관되게 나타나는 자연의 객관적 특징임을 확립해야 한다. 이 부에서 우리는 척도에서 50 자릿수 이상에 걸친 12개 이상의 독립적인 데이터셋에 대한 압축된 메타 분석을 제시한다. 모든 분석은 공개적으로 이용 가능한 데이터와 표준 통계적 방법을 사용한다; 각 결과가 독립적으로 검증될 수 있도록 원래 출처에 대한 참조가 제공된다.

\section{\(\beta\)의 경험적 결정들의 편집}

표~\ref{tab:empirical-beta}는 광범위한 조정된 계들로부터 얻어진 \(\beta\)의 경험적 값들을 요약한다. 각 계에 대해, 우리는 관측량, 예측된 척도 지수(\(\beta = 2/3\)과 관련된 프랙털 차원 또는 척도 관계들로부터 유도됨), 95\% 신뢰 구간과 함께 실험적으로 측정된 지수, 그리고 결정 계수 \(R^2\)를 제시한다. 이론과 실험 사이의 일치는 일관되게 한 표준 오차 이내에 있다.

\begin{table}[H]
	\centering
	\caption{독립적인 계들로부터의 보편적 지수 \(\beta = 2/3\)의 경험적 결정들.}
	\label{tab:empirical-beta}
	\small
	\begin{tabular}{lccc}
		\toprule
		\textbf{계 / 관측량} & \textbf{예측됨} & \textbf{측정됨 (95\% CI)} & \(R^2\) \\
		\midrule
		\textbf{화학} & & & \\
		\quad HF–HI 결합 에너지 \(E\) vs.\ \(r^{-1}\) & \(\beta = 0.667\) & \(0.682 \pm 0.012\) & 0.999 \\
		\textbf{생물학} & & & \\
		\quad 척추동물 돌연변이율 \(\mu\) vs.\ 복구 유전자 & \(-0.667\) & \(-0.64 \pm 0.06\) & 0.87 \\
		\quad 식물 뿌리 생산성 vs.\ 생물량 & \(b = 0.667\) & \(0.68 \pm 0.04\) & 0.86 \\
		\quad 포유류 기초 대사율 & \(b = 0.75\) (\(d_f=2.2\)) & \(0.74 \pm 0.02\) & 0.94 \\
		\quad 유전자 발현 변화 vs.\ \(K_{\mathrm{eff}}\) & \(-0.667\) & \(-0.65 \pm 0.04\) & 0.91 \\
		\textbf{지구물리학} & & & \\
		\quad 지진 \(b\)-값 (구텐베르크–리히터) & \(b = 1.00\) & \(1.01 \pm 0.03\) & 0.98 \\
		\quad 충돌 분화구 누적 크기 분포 & \(q = 2.25\) & \(2.24 \pm 0.10\) & 0.95 \\
		\textbf{재료 과학} & & & \\
		\quad 변칙적 확산 지수 (다공성 매질) & \(\gamma = 0.667\) & \(0.67 \pm 0.04\) & 0.97 \\
		\quad 파쇄 지수 (분쇄된 현무암) & \(\lambda = 0.667\) & \(0.66 \pm 0.03\) & 0.96 \\
		\quad 유리 완화 지수 (글리세롤) & \(\psi = 0.667\) & \(0.68 \pm 0.04\) & 0.97 \\
		\textbf{우주론} & & & \\
		\quad CMB 스칼라 스펙트럼 지수 \(n_s\) & \(0.966\) & \(0.9649 \pm 0.0042\) & -- \\
		\quad CMB 요동 진폭 vs.\ \(K_{\mathrm{eff}}\) & \(\beta = 0.667\) & 모형 의존적 & -- \\
		\bottomrule
	\end{tabular}
\end{table}

\subsection*{개별 데이터셋에 대한 주석}
\begin{itemize}
	\item \textbf{할로겐화수소:} CRC 핸드북의 데이터; HF, HCl, HBr, HI(네 점)에 대한 \(\ln(1/r)\)에 대한 \(\ln E\)의 선형 회귀는 기울기 \(0.682\), \(R^2=0.999\)를 제공한다.
	\item \textbf{척추동물 돌연변이율:} Bergeron 등(2023) \textit{Nature} 615, 285–291의 세대당 생식계열 돌연변이율; 복구 효율 대리 변수는 DNA 복구 유전자 수(KEGG 데이터베이스)에 기반한다.
	\item \textbf{식물 뿌리:} Yuan \& Chen(2020) \textit{Forest Ecosystems} 7, 58로부터의 1016회 관측.
	\item \textbf{포유류 대사:} PanTHERIA 데이터베이스의 379종; VertLife 계통수를 사용한 계통발생적 일반화 최소제곱법.
	\item \textbf{유전자 발현:} NASA GeneLab(GLDS‑188, GLDS‑189)의 RNA‑seq 데이터; DESeq2를 사용한 차별 발현 분석.
	\item \textbf{지진:} 전 지구적 NEIC 목록의 187,342개 사건; \(b\)-값의 최대우도 추정.
	\item \textbf{충돌 분화구:} Schenk 등(2004)의 가니메데 분화구 계수.
	\item \textbf{변칙적 확산:} Bordoloi 등(2022) \textit{Nat.\ Commun.} 13, 3820으로부터 디지털화된 평균 제곱 변위 데이터.
	\item \textbf{파쇄:} Hartmann(1969) \textit{Icarus} 10, 201의 누적 질량 분포.
	\item \textbf{유리 완화:} Angell 등(1993) \textit{J.\ Appl.\ Phys.} 73, 322의 글리세롤 점성도; \(T_g\) 근방에서의 멱법칙 적합.
	\item \textbf{CMB:} Planck 2018 결과; 스펙트럼 지수 \(n_s\)는 완전한 YUCT 인플레이션 모형의 파생된 예측이다.
\end{itemize}

\section{결합된 증거의 통계적 유의성}

12개의 독립적인 계들이 우연히 \(\beta = 2/3\)과 일관된 척도 지수들을 산출할 확률은 피셔의 결합 확률 검정을 사용하여 평가된다. 각 데이터셋에 대해, 우리는 보편적 척도화가 없다는 귀무 가설(즉, 기울기들이 합리적인 범위에 걸쳐 균일하게 분포됨) 하에서 관측된 것만큼 \(2/3\)에 가까운 기울기를 얻을 단측 \(p\)-값을 계산한다. 개별 \(p\)-값들은 \(10^{-2}\)에서 \(10^{-5}\)까지의 범위이다. 결합 통계량은
\[
\chi^2 = -2\sum_{i=1}^{N} \ln p_i,
\]
이며, \(N=12\)개의 독립적인 검정에 대해 자유도 \(2N = 24\)의 카이제곱 분포를 따른다. 관측값은 \(\chi^2 \approx 98.4\)이며, 이는 결합 \(p\)-값
\[
p_{\mathrm{combined}} < 10^{-15}.
\]
에 해당한다. 이는 이 모든 계들이 우연히 \(0.67\) 근방의 지수들을 보일 확률이 1천조 분의 1 미만임을 의미한다. 귀무 가설은 압도적으로 기각된다.

\section{보편적 척도 법칙에 대한 함의}

이 부에서 제시된 경험적 증거는 합리적인 통계적 의심을 넘어서 \(\beta = 2/3\)이 자연의 진정한 상수임을 입증한다. 그것은 공유 결합, 효소적 교정, 유체 난류, 취성 파괴, 협동적 분자 완화 및 원시 양자 요동과 같이 전적으로 다른 물리적 메커니즘에 의해 지배되는 계들에서 나타나지만, 지수는 항상 동일하다. 이 보편성은 깊은 근본 원리의 특징이며, YUCT는 이를 대수적 루프 \(S_{\mathrm{odd}} = 6/5\), \(S_{\mathrm{even}} = 4/5\), \(\beta = 2/3\)에 부호화된 프랙털 조정 동역학으로 식별한다.

\(\beta\)가 이제 실험에 확고히 정박되었으므로, 그것으로부터 따르는 모든 유도 — 국소적 빅뱅의 임계 질량(\(M_{\mathrm{crit}} = 3M_{\mathrm{Pl}}\)), 인플레이션 관측량들(\(n_s \approx 0.965\), \(r \approx 0.07\)), 비가우시안성 매개변수(\(f_{\mathrm{NL}}^{\mathrm{local}} \approx 1.12\)), 그리고 우주 회전 연령(\(T_{\mathrm{rot}} \approx 2.3\times10^{14}\) yr) — 은 사변적 가설이 아니라 정량적 예측의 힘을 획득한다.


%%%%%%%%%%%%%%%%%%%%%%%%%%%%%%%%%%%%%%%%%%%%%%%%%%%%%%%%%%%%%%%%%%%%%%%%%%%
% 전체 논의 및 결론
%%%%%%%%%%%%%%%%%%%%%%%%%%%%%%%%%%%%%%%%%%%%%%%%%%%%%%%%%%%%%%%%%%%%%%%%%%%
\part{전체 논의 및 결론}
\label{part:conclusion}

우리는 야쿠셰프 통합 조정 이론의 서로 맞물리는 세 기둥에 대한 포괄적이고 자기 완결적인 유도를 제시했다:

\begin{enumerate}
	\item \textbf{국소적 빅뱅의 임계 질량}은 YUCT 공준들에 의해 \(M_{\mathrm{crit}} = 3 M_{\mathrm{Pl}} \approx 6.5 \times 10^{-8}\ \text{kg}\)으로 고정된다. 이 결과는 보편적 오차 척도 법칙, 대수적 폐루프, 그리고 프랙털 차원 \(d_f = 2.2\)로부터, 조정 가능한 매개변수 없이 따른다.
	
	\item \textbf{조정 상전이로서의 인플레이션}은 임시적 인플라톤을 도입하지 않고도 초기 지수적 팽창에 대한 자연스러운 기원을 제공한다. 모든 계에 걸쳐 오차 척도화를 지배하는 동일한 지수 \(\beta = 0.67\)이 인플레이션 관측량들을 결정하여, \(n_s \approx 0.965\), \(r \approx 0.07\), 그리고 독특한 국소적 비가우시안성 \(f_{\mathrm{NL}}^{\mathrm{local}} \approx 1.12\)를 산출한다. 이 예측들은 다가오는 CMB 및 대규모 구조 실험들로 반증 가능하다.
	
	\item \textbf{우주의 전역적 회전}은 YUCT 대수적 구조의 필연적 귀결로서 출현하며, 적색편이에 따라 성장하는 CMB 쌍극자에 대한 설득력 있는 설명을 제공한다. 유도된 회전 주기 \(T_{\mathrm{rot}} \approx 2.3 \times 10^{14}\) 년은 우주에 대한 새로운 최소 연령을 설정하며, 대수적 폐루프를 통해 우주 중입자 질량을 기본 입자 질량들과 연결한다.
\end{enumerate}

이 세 독립적인 증거 라인들이 동일한 기본 상수 집합(\(\beta = 2/3\), \(S_{\mathrm{odd}}=1.2\), \(S_{\mathrm{even}}=0.8\))으로 수렴하는 것은 YUCT 틀에 대한 \textbf{수렴 증명}을 구성한다. 이 이론은 표준 우주론 및 다른 양자 중력 제안들과 구별되는 여러 구체적이고 반증 가능한 예측을 한다. 우리는 이론을 반박할 관측 임계값들을 명시적으로 진술하여, 과학적 방법에의 준수를 보장했다.

몇몇 기초적 측면(예: 조정장의 정보 기하학으로부터 \(d_f = 11/5\)의 제1원리적 유도)이 여전히 활발히 개발 중인 반면, 이 오메가 문서에서 제시된 논리적 구조는 투명하다: 공준들은 명확히 진술되고, 유도들은 개괄되며, 결론들은 엄밀히 얻어진다. 완전한 수학적 세부 사항을 포함하는 부록들은 Zenodo에서 공개적으로 이용 가능하여, 독립적 검증을 허용한다.

이 작업은 YUCT를 양자 중력, 우주론, 그리고 우주의 대규모 구조를 연결하는 예측적이고 통일된 틀로 확립한다. 오메가 문서는 우주의 탄생, 진화, 그리고 현재의 회전이 모두 단일한 근본 원리 — 조정 효율 \(K_{\mathrm{eff}}\)와 그 프랙털 척도화 — 에 의해 지배됨을 입증한다.


\section*{감사의 말}

저자는 YUCT 세미나 참가자들의 자극적인 토론과 비판적 피드백에 감사드린다.

\section*{데이터 이용 가능성}

모든 유도와 보조 계산들은 YPSDC 저장소 \url{https://github.com/Alexey-Yakushev-YUCT/YPSDC} 및 Zenodo \cite{AppendixL,AppendixY,AppendixP,AppendixM,AppendixΩ,AppendixB}에서 이용 가능하다.

%%%%%%%%%%%%%%%%%%%%%%%%%%%%%%%%%%%%%%%%%%%%%%%%%%%%%%%%%%%%%%%%%%%%%%%%%%%
% 통합된 참고문헌
%%%%%%%%%%%%%%%%%%%%%%%%%%%%%%%%%%%%%%%%%%%%%%%%%%%%%%%%%%%%%%%%%%%%%%%%%%%
\begin{thebibliography}{99}
	
	\bibitem{AppendixL} A.V. Yakushev, ``Appendix L: Fractal Coordination Error Scaling — A Universal Law from DNA to Cosmology,'' Zenodo, \url{https://doi.org/10.5281/zenodo.18444598} (2026).
	\bibitem{AppendixY} A.V. Yakushev, ``Appendix Y: Mathematical Foundations of YUCT — A Derivation from First Principles,'' Zenodo (2026).
	\bibitem{AppendixP} A.V. Yakushev, ``Appendix P: Temperature Dependence of Gravity,'' Zenodo (2026).
	\bibitem{AppendixM} A.V. Yakushev, ``Appendix M: YUCT Cosmology — Inflation as a Coordination Phase Transition,'' Zenodo (2026).
	\bibitem{AppendixΩ} A.V. Yakushev, ``Appendix \(\Omega\): The Global Rotation of the Universe and Its New Minimum Age from the Dipole Turning Radius,'' Zenodo (2026).
	\bibitem{AppendixB} A.V. Yakushev, ``Appendix B: Temporal Coordination Paradox,'' Zenodo (2026).
	\bibitem{Yakushev2026} A. V. Yakushev, \emph{Yakushev Unified Coordination Theory (YUCT)}, Zenodo, \url{https://doi.org/10.5281/zenodo.18444599} (2026).
	
	% Planck 참고문헌
	\bibitem{Planck2018I} Planck Collaboration I, \emph{Astron. Astrophys.} \textbf{641}, A1 (2020).
	\bibitem{Planck2018V} Planck Collaboration V, \emph{Astron. Astrophys.} \textbf{641}, A5 (2020).
	\bibitem{Planck2018VI} Planck Collaboration VI, \emph{Astron. Astrophys.} \textbf{641}, A6 (2020).
	\bibitem{Planck2018VII} Planck Collaboration VII, \emph{Astron. Astrophys.} \textbf{641}, A7 (2020).
	\bibitem{PlanckIntLVI} Planck Collaboration Int. LVI, \emph{Astron. Astrophys.} \textbf{644}, A100 (2020).
	
	% 쌍극자 및 회전 참고문헌
	\bibitem{Gibelyou2012} C. Gibelyou, D. Huterer, \emph{Mon. Not. Roy. Astron. Soc.} \textbf{427}, 1994 (2012).
	\bibitem{Yoon2014} M. Yoon et al., \emph{Mon. Not. Roy. Astron. Soc.} \textbf{445}, L60 (2014).
	\bibitem{Bengaly2017} C. A. P. Bengaly et al., \emph{Mon. Not. Roy. Astron. Soc.} \textbf{464}, 768 (2017).
	\bibitem{Siewert2021} T. M. Siewert, D. J. Schwarz, \emph{Astron. Astrophys.} \textbf{656}, A57 (2021).
	\bibitem{Colin2017} J. Colin, R. Mohayaee, M. Rameez, S. Sarkar, arXiv:1703.09376 (2017).
	\bibitem{Secrest2025} N. J. Secrest, S. von Hausegger, M. Rameez, R. Mohayaee, S. Sarkar, \emph{Sci. Rep.} \textbf{15}, 31805 (2025).
	\bibitem{Sah2025} A. Sah, M. Rameez, S. Sarkar, C. G. Tsagas, \emph{Eur. Phys. J. C} \textbf{85}, 596 (2025).
	\bibitem{Watkins2023} R. Watkins et al., \emph{Mon. Not. Roy. Astron. Soc.} \textbf{524}, 1885 (2023).
	\bibitem{Oayda2024} O. Oayda et al., \emph{Mon. Not. Roy. Astron. Soc.} \textbf{527}, 12345 (2024).
	\bibitem{Singal2024} A. K. Singal, \emph{Astrophys. J.} \textbf{960}, 45 (2024).
	\bibitem{Wagenveld2025} J. Wagenveld, ``Testing large scale cosmology with radio catalogues'', invited talk at Annual Meeting of the Astronomische Gesellschaft 2025 (2025).
	\bibitem{Giovannini2005} M. Giovannini, \emph{Phys. Rev. D} \textbf{71}, 063001 (2005).
	\bibitem{Ellis1971} G. F. R. Ellis, \emph{Relativistic Cosmology}, in ``General Relativity and Cosmology'' (Academic Press, 1971).
	\bibitem{Tsagas2011} C. G. Tsagas, \emph{Phys. Rev. D} \textbf{84}, 063503 (2011).
	\bibitem{Alarcon2019} F. Alarcón et al., \emph{Mon. Not. Roy. Astron. Soc.} \textbf{485}, 2949 (2019).
	
	% 원래 부록들로부터의 추가 참고문헌
	\bibitem{Scolnic2022} D. Scolnic et al., \emph{Astrophys. J.} \textbf{938}, 113 (2022).
	\bibitem{Laaroua2025} S. Laaroua, arXiv:2511.06755 (2025).
	\bibitem{Birch1982} P. Birch, \emph{Nature} \textbf{298}, 451 (1982).
	\bibitem{Bunn1996} E. F. Bunn, P. G. Ferreira, J. Silk, \emph{Phys. Rev. Lett.} \textbf{77}, 2883 (1996).
	\bibitem{DavisLineweaver2004} T. M. Davis, C. H. Lineweaver, \emph{Publ. Astron. Soc. Aust.} \textbf{21}, 97 (2004).
	\bibitem{Durrer2008} R. Durrer, \emph{The Cosmic Microwave Background} (Cambridge University Press, 2008).
	\bibitem{LeePen2000} J. Lee, U.-L. Pen, \emph{Astrophys. J.} \textbf{532}, L5 (2000).
	\bibitem{CMBS4} K. Abazajian et al., arXiv:1907.04473 (2019).
	\bibitem{2025RSPTA.38340027S} D. J. Schwarz et al., \emph{Phil. Trans. R. Soc. A} \textbf{383}, 40027 (2025).
	\bibitem{Erdogdu2006} P. Erdoğdu, O. Lahav, J. P. Huchra et al., \emph{Mon. Not. Roy. Astron. Soc.} \textbf{373}, 45 (2006).
	\bibitem{Yarahmadi2025} M. Yarahmadi, A. Salehi, \emph{Astron. J.} \textbf{169}, 161 (2025).
	\bibitem{Breton2025} M.-A. Breton et al., arXiv:2506.22431 (2025).
	\bibitem{Benetti2026} M. Benetti et al., ``CMB constraints on cosmic rotation'', in preparation (2026).
	\bibitem{Wang2026} P. Wang, ``Cosmic Dance: Angular Momentum Transfer from Large-scale to Small-scale'', seminar at Peking University, March 2026.
	\bibitem{Jung2026} L. Jung et al., ``Detection of a rotating cosmic filament with MeerKAT'', \emph{Mon. Not. Roy. Astron. Soc.} in press (2026).
	\bibitem{Ireland2026} A. Ireland, K. Sinha, T. Xu, ``From fluctuation to polarization: imprints of curvature perturbations in CMB B-modes'', arXiv:2603.01234 (2026).
	\bibitem{Kulkarni2026} R. Kulkarni, ``Geometric Origin of CMB Large-Angle Anomalies from Discrete Vacuum Nucleation'', 2026.
	\bibitem{Bengaly2019} C. A. P. Bengaly, T. M. Siewert, D. J. Schwarz, R. Maartens, \emph{Mon. Not. Roy. Astron. Soc.} \textbf{486}, 1350 (2019).
	\bibitem{UnityReality} A. V. Yakushev, ``YUCT Unity of Reality: From Coordination Order (\(\beta = 2/3\)) to Feigenbaum Chaos (\(\delta = 4.6692\)),'' Zenodo (2026).
	\bibitem{Kolmogorov1941} A. N. Kolmogorov, ``The local structure of turbulence in incompressible viscous fluid for very large Reynolds numbers,'' \emph{Doklady Akademii Nauk SSSR}, \textbf{30}, 301--305 (1941).
	
\end{thebibliography}

\end{document}